\documentclass[11pt]{article}
% LaTeX-only set-up (the web build does not read this file).
\usepackage[T1]{fontenc}
\usepackage[utf8]{inputenc}
\usepackage{lmodern}
\usepackage{amsmath,amssymb,amsthm}
\usepackage[a4paper,margin=1.1in]{geometry}
\usepackage{microtype}
\usepackage{booktabs}
\usepackage{enumitem}
\usepackage{newunicodechar}
\usepackage[hidelinks]{hyperref}
\setlist{itemsep=0.2em,topsep=0.4em}
% Statements carry their own explicit label through \head, so nothing is numbered automatically.
\newenvironment{statement}{\par\addvspace{\medskipamount}\noindent\ignorespaces}{\par\addvspace{\medskipamount}}
\newenvironment{frontnote}{\par\addvspace{\medskipamount}\small\noindent\ignorespaces}{\par\addvspace{\medskipamount}}
\renewcommand{\qedsymbol}{$\blacksquare$}
% Unicode that may appear inside verbatim and \texttt (Lean code)
\newunicodechar{α}{\ensuremath{\alpha}}\newunicodechar{β}{\ensuremath{\beta}}
\newunicodechar{γ}{\ensuremath{\gamma}}\newunicodechar{δ}{\ensuremath{\delta}}
\newunicodechar{μ}{\ensuremath{\mu}}\newunicodechar{ω}{\ensuremath{\omega}}
\newunicodechar{τ}{\ensuremath{\tau}}\newunicodechar{λ}{\ensuremath{\lambda}}
\newunicodechar{Π}{\ensuremath{\Pi}}\newunicodechar{Σ}{\ensuremath{\Sigma}}
\newunicodechar{∀}{\ensuremath{\forall}}\newunicodechar{∃}{\ensuremath{\exists}}
\newunicodechar{∧}{\ensuremath{\wedge}}\newunicodechar{∨}{\ensuremath{\vee}}
\newunicodechar{¬}{\ensuremath{\neg}}\newunicodechar{→}{\ensuremath{\to}}
\newunicodechar{↔}{\ensuremath{\leftrightarrow}}\newunicodechar{≤}{\ensuremath{\le}}
\newunicodechar{∈}{\ensuremath{\in}}\newunicodechar{⊢}{\ensuremath{\vdash}}
\newunicodechar{₂}{\ensuremath{_2}}\newunicodechar{₃}{\ensuremath{_3}}\newunicodechar{⁺}{\ensuremath{^+}}
% Lean listings are wider than the text font allows, so set them smaller.
\makeatletter\def\verbatim@font{\normalfont\ttfamily\footnotesize}\makeatother

% Shared by the PDF and the web build. Keep to simple \newcommand definitions.
% The short name of the main notion (alpha-plus-recursively Mahlo). Change it here and nowhere else.
% It is used both as an adjective ("the least \RM{} ordinal") and as a noun at the start of
% sentences, so it has to be an abbreviation or symbol, not a lower-case word.
\newcommand{\RM}{RM\textsuperscript{+}}
% Bold run-in head of a statement or front note: \head{Lemma 4.1}
\newcommand{\head}[1]{\textbf{#1.}}
% Named classes and predicates
\newcommand{\Mh}{\mathrm{Mh}}
\newcommand{\Sat}{\mathrm{Sat}}
\newcommand{\Mem}{\mathrm{Mem}}

% Characters of the Lean listings in the appendix not covered by the shared preamble (PDF only).
\newunicodechar{η}{\ensuremath{\eta}}\newunicodechar{θ}{\ensuremath{\theta}}
\newunicodechar{ξ}{\ensuremath{\xi}}\newunicodechar{ζ}{\ensuremath{\zeta}}
\newunicodechar{∣}{\ensuremath{\mid}}\newunicodechar{≠}{\ensuremath{\neq}}
\title{The least ordinal that is $\Pi_n$-reflecting of degree $\delta_\alpha$ is the least $\Pi_{n+1}$-reflecting ordinal}
% The date is stated here once; the PDF and the web page both take it from this line.
\newcommand{\paperversion}{2026-11-22}
\date{\paperversion}
\begin{document}
% Local notation (kept here so that the web build, which reads only the body, sees it too).
\newcommand{\Ad}{\mathrm{Ad}}
\newcommand{\RMc}{\Mh(2)}
\newcommand{\Lim}{\mathrm{Lim}\,}
\newcommand{\CA}{C_{\mathrm{Ad}}}
\newcommand{\Lv}{\mathcal{L}}
\newcommand{\Cert}{\mathrm{Cert}}
\newcommand{\Tab}{\mathrm{Tab}}
\maketitle

\begin{frontnote}\head{Abstract} For a $\Sigma_1$ presentation $P$ with parameters, a node $a$ and $n \ge 1$, let an
ordinal $\mu$ be in the class $C_n(P,a)$ if it is $\Pi_n$-reflecting on the classes of the predecessors of $a$, the
relation of $P$ being evaluated afresh in $L_\mu$. We show that for every $n \ge 1$ the least limit ordinal $\alpha >
\omega$ that lies in the class of every node in the externally well-founded part of every such presentation at
$\alpha$ is the least $\Pi_{n+1}$-reflecting ordinal. The case $n = 2$ is an earlier paper. The step beyond it is the
lower bound for $n \ge 3$: a search tree whose promises are kept by choosing witnesses, each choice removing two
quantifier blocks, so that only a bounded number of steps must be unwound. For admissible $\alpha > \omega$ the ranks
of these nodes are exactly the ordinals below $\delta_\alpha$, the supremum of the order types of $\Delta_1(L_\alpha)$
well-orderings definable with parameters, and we call the condition $\Pi_n$-reflection of degree $\delta_\alpha$. At
admissible $\gamma > \omega$ in which every earlier level is
countable, $\gamma$ is $\Pi_n$-reflecting of degree $\delta_\gamma$ exactly when it is $\Pi_{n+1}$-reflecting or
$\Pi_n$-reflecting on the $\Pi_{n+1}$-reflecting ordinals ($n \ge 2$). With ``admissible limit of'' in place of
``$\Pi_n$-reflecting on'', the least ordinal of degree $\delta_\alpha$ is the least recursively Mahlo ordinal.
\end{frontnote}

\begin{frontnote}\head{Provenance} The proofs were found and written by AI systems. No human specialist has
reviewed them, and novelty is not established. Section 7 states what has been checked and how.
\end{frontnote}

\section{Introduction}

\textbf{Main theorem} (Theorems 4.3, 4.8 and 4.18). \emph{For every $n \ge 1$, the least limit ordinal $\alpha >
\omega$ such that
$\alpha \in C_n(P, a)$ for every $\Sigma_1$ presentation $P$ with parameter in $L_\alpha$ and every node $a \in L_\alpha$
in the well-founded part of $P$ at $\alpha$ is the least $\Pi_{n+1}$-reflecting ordinal $K_n$.} Here a presentation is
a formula with a parameter, read as a predecessor relation on $L_\mu$ at each level $\mu$; well-foundedness is taken in
the universe; and $\mu \in C_n(P,a)$ when $\mu$ is $\Pi_n$-reflecting on $C_n(P,c)$ for every predecessor $c$ of $a$
in $L_\mu$, the relation being evaluated afresh at each level (Definition 2.1). For $n = 2$ this is the main theorem of
the earlier paper $[E]$, \emph{The least $\alpha^+$-recursively Mahlo ordinal is the least $\Pi_3$-reflecting
ordinal}, cited as [E, Theorem B]; we use its conventions and cite its items as [E, Lemma 4.1].
That $K_n$ has the property is the known upper-bound argument (Theorem 4.3). The step beyond $[E]$ is the lower bound
for $n \ge 3$ (Section 4.3): the search tree of $[E]$ must now keep promises with more than two quantifier blocks, and does so by
nesting them. For $n = 1$ the lower bound is elementary (Section 4.2).

For admissible $\alpha > \omega$ the ranks of the nodes in the condition are exactly the ordinals below
$\delta_\alpha$, the supremum of the order types of $\Delta_1(L_\alpha)$ well-orderings (Theorem 3.1), and we call the
condition \emph{$\Pi_n$-reflection of degree $\delta_\alpha$}. The name records only which ranks are available at
$\alpha$: the class of a node is not determined by its rank at $\alpha$, because it depends on how the relation is
interpreted at the smaller levels (Examples 4.7 and 5.10). When $\alpha$ is countable, locally countable and below the
least $\Sigma^1_1$-reflecting ordinal, $\delta_\alpha = \alpha^+$ (Theorem 3.2), which is why $[E]$ calls the case $n =
2$ ``$\alpha^+$-recursively Mahlo''.

\textbf{Characterisation} (Theorem 4.21). \emph{Let $n \ge 2$, and let $\gamma > \omega$ be admissible with every
$L_\eta$, $\eta < \gamma$, countable in $L_\gamma$. Then $\gamma$ is $\Pi_n$-reflecting of degree $\delta_\gamma$ if
and only if $\gamma$ is $\Pi_{n+1}$-reflecting or $\Pi_n$-reflecting on the $\Pi_{n+1}$-reflecting ordinals.} The
countability hypothesis holds at every admissible $\gamma$ up to the least ordinal that is $\Pi_m$-reflecting for every
$m$, and fails at large admissible $\gamma$. For $n = 2$ this is [E, Theorem C].

\textbf{Variants.} Each operation on classes below has a \emph{next level}: for ``limit of'', which is
$\Pi_1$-reflection on a class, it is $\Pi_2$-reflection, that is, admissibility; for ``$\Pi_n$-reflecting on'' it is
$\Pi_{n+1}$-reflection, so that for this family the least point advances from $\Pi_n$ to $\Pi_{n+1}$; for ``admissible
limit of'' it is $\Pi_2$-reflection on the admissible ordinals, that is, recursive Mahloness. In each case the least
ordinal of degree $\delta_\alpha$ is the least ordinal of the next level. In each stated range the ordinals of that
degree form the class $B \cup O(B)$, where $B$ is the class of the next level and $O(B)$ is the result of applying the
operation to $B$; the exceptions, the members not in $B$, are those of $O(B) \setminus B$.
\begin{itemize}
\item \emph{Limit of} ($n = 1$). Least element $\omega_1^{CK}$ (Theorem 4.8). For limits $\gamma$ such that $L_\gamma$
  contains a well-ordering of $\omega$ of type at least the supremum of $\omega$ and the admissible ordinals below
  $\gamma$: of
  degree $\delta_\gamma$ if and only if admissible or a limit of admissibles (Theorem 4.9). Among all limits $\gamma$
  of degree $\delta_\gamma$, the least non-admissible one is $\omega_\omega^{CK}$ (Corollary 4.11).
\item \emph{$\Pi_n$-reflecting on}, $n \ge 2$. Least element $K_n$ (Theorem 4.18); characterisation above. The least
  member that is not $\Pi_{n+1}$-reflecting is the least ordinal $\gamma_1(n)$ that is $\Pi_n$-reflecting on the
  $\Pi_{n+1}$-reflecting ordinals.
\item \emph{Admissible limit of} (recursive inaccessibility). Least element $M$, the least recursively Mahlo ordinal
  (Theorem 5.8). For admissible $\gamma$ up to the least ordinal $\Pi_m$-reflecting for every $m$: of degree
  $\delta_\gamma$ if and only if recursively Mahlo or a limit of recursively Mahlo ordinals (Corollary 5.9). The least
  member that is not recursively Mahlo is the least admissible limit of recursively Mahlo ordinals.
\end{itemize}
At admissible $\alpha$, restricting to $\Sigma_1$ linear orderings gives the same notion of $\Pi_n$-reflection of
degree $\delta_\alpha$ (Proposition 4.23).

\textbf{How the proofs go.} For the upper bounds, membership in the class of a node is expressed by a sentence, and
reflecting that sentence gives the class of every predecessor, by induction on rank. Theorem 4.3 reflects the
$\Pi_{n+1}$ membership formula of Lemma 4.2. Theorem 5.7(a) reflects a $\Pi_2$ membership formula onto the admissible
ordinals, with a given bound named as a constant. For the lower bounds, a $\Sigma_1$ tree searches for a reflecting
level, an ordinal of the next level, at or below the level under test, and a descent shows that the level under test is
not in the class of its root. For $n \ge 2$ the descent passes a $\Pi_n$
sentence down (Section 4.3); for ``limit of'' a set relation of large rank suffices (Section 4.2); for ``admissible
limit of'' the descent carries one set down to an admissible target (Section 5). Appendix A lists which result uses
which.

\textbf{Nearest results.} Arai (Iterating the recursively Mahlo operations, Definition 1.1 and Lemma 1.2) iterates the
recursively Mahlo operation along a relation evaluated in the current universe, and shows that a $\Pi_3$-reflecting
universe lies in the class of every $n \in \omega$ up to which transfinite induction along a $\Sigma_1$ relation on $\omega$
holds; his Theorem 1.3 reduces the theory of $\Pi_3$-reflecting universes proof-theoretically to limits of admissibles
with such iterations along one $\Sigma_1$ relation. Theorem 4.3 is that argument for every $n$, for presentations with
parameters and with the well-founded part taken externally; Arai's Lectures (Lemma 3.33 and Remark 3.34) and later
papers (Conservations, Theorem 2.2; Mahlo classes, Section 1) give inclusions of the same kind, with degrees given by
ordinals or notation codes. Richter and Aczel (Theorems 1.9--1.11) characterise $\Pi_1$- and $\Pi_2$-reflection on a class and show
that recursively Mahlo ordinals lie in the iterates of ``admissible limit of''; this is the upper bound of Section 5 in
another form. Rathjen (Proof theory of reflection, 1993 preprint) shows that every $\Pi_3$-reflecting $\kappa$ is
recursively $\kappa$-Mahlo (Section 2) and, for a weakly compact cardinal $K$, that the Mahlo classes of degree below
$\varepsilon_{K+1}$, with degrees evaluated through hulls, are stationary (Theorem 4.12). All of these are upper bounds
or proof-theoretic counterparts. We did not find in them a lower bound of the kind proved here, the characterisations,
or the rank statement of Theorem 3.1; Aguilera states without proof that $\delta_\alpha$ is unchanged if
$\alpha$-recursively enumerable well-orderings are allowed.

\section{Preliminaries}

\subsection{\texorpdfstring{Conventions taken from $[E]$}{Conventions taken from [E]}}

[E, Section 2] is used without change. This section adds only the truth predicates $\Sat_m$ for every $m$, built as
in [E, (C4)], and the reading of operations on codes from the truth table at $\omega$. In brief:

\begin{itemize}
\item $L_\alpha$ is the $\alpha$-th constructible level. A set is \emph{countable in} $L_\gamma$ if it is empty or the
  image of a surjection from $\omega$ that belongs to $L_\gamma$. A sentence is evaluated in $L_\beta$ only if its
  constants are elements of $L_\beta$.
\item \emph{Formula classes.} A formula is $\Pi_n$ if it is logically equivalent to one with $n$ alternating blocks
  of unbounded quantifiers, the first universal, followed by a $\Delta_0$ matrix. In a limit level a block of like
  quantifiers can be contracted to one quantifier over tuples and a list of constants to one tuple, so at limit levels
  we write every $\Pi_m$ sentence as $\forall x_1\, \exists x_2 \cdots x_m\, \theta(\vec x, q)$, with one variable per
  block and one constant, and a $\Pi_j$ sentence with $j < m$ counts as $\Pi_m$ by vacuous quantifiers. Block
  contraction is used at limit levels.
\item \emph{Certificates.} $\Cert(W, \xi)$ is a $\Delta_0$ formula; a \emph{certificate through} $\xi$ is a set $W$
  with $\Cert(W,\xi)$. It consists of syntax data and, for each $\nu \le \xi$, a set $\Lv_\nu$ and a truth table
  $\mathrm{Tr}_\nu$; $[E]$ writes $M_\nu$ for $\Lv_\nu$, and we keep $M$ for an ordinal. (C1): every certificate is the
  true one: $\Lv_\nu = L_\nu$ and $\mathrm{Tr}_\nu$ is the satisfaction relation of $L_\nu$. (C2): an admissible
  $L_\alpha$, $\alpha > \omega$, contains certificates through every $\xi < \alpha$. (C3): $G$ is the $\Pi_2$ sentence
  ``$\omega$ exists, every ordinal has a successor, and every ordinal has a certificate''; it holds at admissible levels
  above $\omega$, and a nonempty level satisfying it is a limit level above $\omega$ (parameter-free tests are given the
  dummy constant $\emptyset$). (C4): $[E]$ gives $\Pi_2$ and $\Pi_3$ truth predicates $\Sat_2$, $\Sat_3$ for levels
  satisfying $G$; we use $\Sat_m$ for every $m$, defined below. (C5): $\Sigma_1$ sentences persist from a level to every
  larger one. (C6): condensation.
\item \emph{Order certificates} [E, Section 8]: the canonical well-ordering $<_L$ is $\Delta_1$ over every admissible
  level above $\omega$.
\item \emph{Countability} [E, Lemma 7.2]: if $\gamma > \omega$ is admissible and some $L_\eta$ with $\eta < \gamma$ is
  not countable in $L_\gamma$, then some ordinal below $\gamma$ is $\Pi_m$-reflecting for every $m$.
\end{itemize}

\textbf{Operations on codes.} Formula codes are hereditarily finite sets. Every operation on codes used below is
primitive recursive: recognising the codes of $\Pi_m$ formulas in the form above and their matrices, the operations
$\mathrm{dual}$ and $\mathrm{tail}$ of Definition 4.13, and the enumerations of tasks in Definitions 4.13 and 5.3. Each
such relation is defined over $L_\omega$ by a fixed formula, and ``the relation holds of given codes'' is the bounded
condition that the pair of that formula and the codes lies in $\mathrm{Tr}_\omega$. So these operations are read from
the truth table at $\omega$ of any certificate through an ordinal $\ge \omega$, exactly as $[E]$ reads its enumeration
of tasks (end of [E, Section 6.1]); no data are added to certificates.

\textbf{Truth predicates.} For $m \ge 1$, $\Sat_m(\varphi, q)$ is $\forall x_1\, \exists x_2 \cdots x_m\, Q\nu\,
QW\, B$, where $Q$ is the last of the $m$ quantifiers and $B$ is a bounded condition. With $\vec x = (x_1, \dots,
x_m)$: if $Q = \exists$,
\[ B = \Cert(W,\nu) \wedge \omega \le \nu \wedge \vec x, q \in \Lv_\nu \wedge \mathrm{Val}_m, \]
and if $Q = \forall$,
\[ B = \neg\Cert(W,\nu) \vee \nu < \omega \vee \neg(\vec x, q \in \Lv_\nu) \vee \mathrm{Val}_m. \]
Here
$\mathrm{Val}_m$ says that, according to $\mathrm{Tr}_\omega$, $\varphi$ is the code of a $\Pi_m$ formula in the form
above with some matrix $\theta \in \Lv_\omega$, and $(\theta, \langle x_1, \dots, x_m, q\rangle) \in \mathrm{Tr}_\nu$.
The quantifiers $\nu, W$ join the last block, so $\Sat_m$ is $\Pi_m$. As in the proof of [E, (C4)]: if $L_\beta \models
G$ and $q \in L_\beta$, then $L_\beta \models \Sat_m(\varphi, q)$ if and only if $\varphi$ is a $\Pi_m$ code in the
form above and $L_\beta \models \varphi(q)$. Indeed, given $x_1, \dots, x_m \in L_\beta$, they and $q$ lie in some
$L_\nu$ with $\omega \le \nu < \beta$, a certificate through $\nu$ is in $L_\beta$ by (C3), every certificate is the true
one by (C1), and $\theta(\vec x, q)$ holds in $L_\nu$ if and only if it holds in $L_\beta$.

From a certificate through $\nu$ one also reads off, by a $\Delta_0$ condition, whether $\nu$ is admissible,
recursively Mahlo or $\Pi_{n+1}$-reflecting, using the truth tables of $L_\nu$ and of smaller levels as in [E, Lemma
7.5]. For such a class $X$, the sentence ``$X$ is unbounded'' means the $\Pi_2$ sentence ``above every ordinal there is
a certificate showing that some $\nu$ is in $X$''; at a level satisfying $G$ it means what it says, by (C1) and (C3).

\subsection{Reflection and classes of ordinals}

An ordinal $\alpha$ is \emph{$\Pi_n$-reflecting} if every $\Pi_n$ sentence $\varphi(q)$ with $q \in L_\alpha$ true in
$L_\alpha$ is true in some $L_\beta$ with $\beta < \alpha$ and $q \in L_\beta$. For a class $X$ of ordinals, $\mu$ is
\emph{$\Pi_n$-reflecting on $X$} if the same holds with $\beta \in X \cap \mu$. As in $[E]$, ordinals called
$\Pi_n$-reflecting or $\Pi_n$-reflecting on a class are limit ordinals above $\omega$, and reflection targets need not
be admissible. A limit $\alpha > \omega$ is $\Pi_2$-reflecting if and only if it is admissible (Richter--Aczel, Theorem
1.8). For any ordinal $\mu$, $\mu^+$ is the least admissible ordinal above $\mu$. An ordinal $\alpha$ is \emph{locally
countable} if every element of $L_\alpha$ is countable in $L_\alpha$.

\begin{itemize}
\item A limit ordinal is nonzero. $\Lim X$ is the class of limit ordinals $\mu$ such that $X \cap \mu$ is unbounded in
  $\mu$; its members are \emph{limits of $X$}.
\item $\Ad$ is the class of admissible ordinals above $\omega$. An ordinal is \emph{recursively inaccessible} if it is
  in $\Ad \cap \Lim \Ad$, that is, an admissible limit of admissible ordinals.
\item $\Mh(\xi)$ is as in [E, Section 3.1]: $\Mh(0)$ is the class of limits above $\omega$, $\Mh(\xi+1)$ the class of
  ordinals $\Pi_2$-reflecting on $\Mh(\xi)$, and $\Mh(\lambda) = \bigcap_{\xi<\lambda}\Mh(\xi)$ at limits. So $\Mh(1) =
  \Ad$, and $\Mh(2)$ is the class of \emph{recursively Mahlo} ordinals, those $\Pi_2$-reflecting on $\Ad$. The members
  of $\Mh(3)$, the ordinals $\Pi_2$-reflecting on the recursively Mahlo ordinals, are called \emph{recursively
  1-Mahlo}, just as ``1-inaccessible'' is one step above ``inaccessible''.
\item $\omega_k^{CK}$, for $1 \le k < \omega$, is the $k$-th admissible ordinal above $\omega$, $\omega_0^{CK} =
  \omega$, and $\omega_\omega^{CK}$ is the least limit of admissible ordinals.
\item $A_n$ is the class of $\Pi_{n+1}$-reflecting ordinals, $K_n = \min A_n$, and $\gamma_1(n)$ is the least ordinal
  that is $\Pi_n$-reflecting on $A_n$. By Richter--Aczel 1.8, $A_1 = \Ad$. $[E]$'s $A$, $K$, $\gamma_1$ are $A_2$,
  $K_2$, $\gamma_1(2)$.
\item A \emph{$\Sigma^1_1$ sentence over $L_\alpha$} is a sentence $\exists S\, \varphi(S, q)$ in which $S$ ranges over
  subsets of $L_\alpha$, $\varphi$ is first-order with a predicate for $S$, and $q \in L_\alpha$; $\alpha$ is
  \emph{$\Sigma^1_1$-reflecting} if every such sentence true in $L_\alpha$ is true in some $L_\beta$, $\beta < \alpha$,
  with $q \in L_\beta$ (Richter--Aczel, Definition 1.7). $\sigma^1_1$ denotes Aguilera's least $\Sigma^1_1$-reflecting
  ordinal, in his countable, locally countable setting (Aguilera, Sections 1--2). It is admissible.
\end{itemize}

\begin{center}\begin{tabular}{p{0.14\linewidth}p{0.76\linewidth}}
\toprule
Ordinal & Definition \\
\midrule
$\omega_1^{CK}$ & the least admissible ordinal above $\omega$, that is, $K_1$ \\
$\omega_\omega^{CK}$ & the least limit of admissible ordinals \\
$M$ & the least recursively Mahlo ordinal \\
$N$ & the least recursively 1-Mahlo ordinal \\
$K_n$, $\gamma_1(n)$ & the least $\Pi_{n+1}$-reflecting ordinal; the least ordinal $\Pi_n$-reflecting on $A_n$ \\
$\xi^*$ & the least ordinal $\Pi_m$-reflecting for every $m$ (with set parameters), equivalently the least $\alpha$ with $L_\alpha \prec_{\Sigma_1} L_{\alpha+1}$ (Richter--Aczel, Theorem 1.18); it exists by [E, proof of Lemma 7.1] \\
\bottomrule
\end{tabular}\\[0.4em]Table 1. Some ordinals used below.\end{center}

By [E, Lemma 7.2], every $L_\eta$ with $\eta < \gamma$ is countable in $L_\gamma$ for every admissible $\gamma \le
\xi^*$.

\subsection{Presentations and their classes}

Presentations are as in [E, Section 3.2]. A \emph{presentation} is a pair $P = (\rho, p)$ of a formula $\rho(x,y,z)$
and a set $p$; at a limit $\mu > \omega$ with $p \in L_\mu$, $c$ is a \emph{predecessor} of $a$ at $\mu$ if $L_\mu
\models \rho(c, a, p)$, and $(c,a)$ is an \emph{edge}. Elements of $L_\mu$ are \emph{nodes}. $P$ is a $\Sigma_1$
presentation if $\rho$ is $\Sigma_1$. The \emph{well-founded part} at $\mu$ is the set of nodes from which there is no
infinite descending sequence, taken in the universe, and $\mathrm{rank}_\mu(a)$ is the rank of $a$ in the relation at
$\mu$. By (C5), an edge of a $\Sigma_1$ presentation at $\mu$ is an edge at every larger level, so a node in the
well-founded part at $\gamma$ is in the well-founded part at every smaller limit $\mu$ containing it and the parameter,
with $\mathrm{rank}_\mu(a) \le \mathrm{rank}_\gamma(a)$. If $\lhd$ is a relation that is a set, $P_\lhd$ is the
$\Delta_0$ presentation ``$(x,y) \in z$'' with parameter $\lhd$; its edges are the same at every level containing
$\lhd$, and we call it a \emph{set relation}.

For a limit $\alpha > \omega$, $\delta_\alpha$ is the supremum of the order types of well-orderings whose field is a
subclass of $L_\alpha$ and which are $\Delta_1$-definable over $L_\alpha$ with parameters (the notation of Aguilera,
Section 2). By Theorem 3.1 below, for admissible $\alpha > \omega$ these order types are exactly the ranks of nodes in
the well-founded parts of $\Sigma_1$ presentations at $\alpha$; this is why the notions of Definition 2.2 are named
after $\delta_\alpha$.

\begin{statement}\head{Definition 2.1} Let $P$ be a presentation. By recursion on $\mu$, for all nodes at once:
\begin{itemize}
\item for $n \ge 1$, $\mu \in C_n(P, a)$ if and only if $\mu$ is a limit ordinal above $\omega$, $p, a \in L_\mu$,
  and $\mu$ is $\Pi_n$-reflecting on $C_n(P, c)$ for every predecessor $c$ of $a$ at $\mu$;
\item $\mu \in \CA(P, a)$ if and only if $\mu \in \Ad$, $p, a \in L_\mu$, and $\mu \in \Lim \CA(P, c)$ for every
  predecessor $c$ of $a$ at $\mu$.
\end{itemize}
\end{statement}

$C_2$ is $[E]$'s class $C$ [E, Definition 3.1]. A node without predecessors at $\mu$ imposes only the first two
clauses.

\begin{statement}\head{Definition 2.2} Let $\alpha > \omega$ be a limit ordinal. $\alpha$ is \emph{$\Pi_n$-reflecting
of degree $\delta_\alpha$} if $\alpha \in C_n(P, a)$ for every $\Sigma_1$ presentation $P$ with $p \in L_\alpha$ and
every node $a \in L_\alpha$ in the well-founded part of the relation defined by $P$ at $\alpha$. It is
\emph{recursively inaccessible of degree $\delta_\alpha$} if the same holds with $\CA$ in place of $C_n$.
\end{statement}

For $n = 2$ this is $[E]$'s $\alpha^+$-recursively Mahlo. For a set relation and a node of rank $\xi$ the class $C_2$
is $\Mh(\xi)$, restricted to the $\mu$ containing the parameter and the node [E, Example 3.2(1)]. Again for a set
relation, the class $C_n$ of a node of rank 0 is the class of limits above $\omega$ containing the parameter and the
node, and every member of the
class of a node of rank 1 is $\Pi_n$-reflecting; the class $\CA$ of a node of rank 0 is $\Ad$, and that of a node of
rank 1 is $\Ad \cap \Lim \Ad$, the recursively inaccessible ordinals, in both cases restricted to the $\mu$ containing
the parameter and the node. This is why the iteration of ``admissible limit of'' is called recursive inaccessibility of a degree. For other $\Sigma_1$ presentations the class of a node is not determined by its
rank (Examples 4.7 and 5.10).

\textbf{Class certificates.} Membership in these classes is computed inside a level by the method of [E, Lemma 5.1,
Steps 1--2], which we state once in general. Let $\Phi(W, \nu, x, H)$ be a $\Delta_0$ formula, possibly with set
parameters. For a certificate $W$ through $\mu$ and $\nu \le \mu$, $W{\restriction}(\nu+1)$ is $W$ cut down to its
rows at indices $\eta \le \nu$; it is a certificate through $\nu$. A \emph{class certificate for $\Phi$ through $\mu$}
is a pair $(W, \Tab)$ of a certificate $W$ through $\mu$ and a function $\Tab$ on $\mu + 1$ with $\Tab(\nu) \subseteq
\Lv_\nu$ (written $D$ in [E]), such that for every $\nu \le \mu$ and $x \in \Lv_\nu$: $x \in \Tab(\nu)$ if and only if
$\Phi(W{\restriction}(\nu+1), \nu, x, \Tab{\restriction}\nu)$. The restrictions are read as bounded views of the given
graphs, since a transitive set containing a graph need not contain its restriction: a quantifier over a restricted
graph is relativised to the given graph by its index guard, equality with a set is tested by the two bounded
inclusions, and membership in a set by bounded search for an equal element. The pair $(W, \Tab)$ has a fixed finite
tuple shape, and induction on the fixed formula $\Phi$ gives a bounded test on the given data. So being a class certificate
is $\Delta_0$.

\begin{statement}\head{Lemma 2.3} Let $\Phi$ be as above, and let $J$ be a class of pairs $(\nu, x)$ with $x \in
L_\nu$; write $J_\nu = \{x : (\nu, x) \in J\}$. Suppose that, for the true certificate $W$ through $\nu$, $x \in J_\nu$
if and only if $\Phi(W, \nu, x, (J_\eta)_{\eta < \nu})$.
\begin{itemize}
\item[(i)] Every class certificate for $\Phi$ through $\mu$ has $\Tab(\nu) = J_\nu$ for all $\nu \le \mu$.
\item[(ii)] If the parameters of $\Phi$ lie in $L_\alpha$, every admissible $L_\alpha$, $\alpha > \omega$, contains a
  class certificate for $\Phi$ through every $\mu < \alpha$.
\item[(iii)] Thus ``there is a class certificate through $\mu$ with $x \in \Tab(\mu)$'' is a $\Sigma_1$ formula. In any
  transitive set containing its parameters, its truth implies $x \in J_\mu$. Under the hypotheses of (ii), it holds in
  $L_\alpha$ whenever $\mu < \alpha$ and $x \in J_\mu$.
\end{itemize}
\end{statement}

\begin{proof}
(i) By (C1), $W$ is the true certificate, so $W{\restriction}(\nu+1)$ is the true certificate through $\nu$. By
induction on $\nu \le \mu$: if $\Tab(\eta) = J_\eta$ for $\eta < \nu$, then $x \in \Tab(\nu)$ if and only if
$\Phi(W{\restriction}(\nu+1), \nu, x, (J_\eta)_{\eta<\nu})$, that is, $x \in J_\nu$. (ii) By (C2) a certificate $W$
through $\mu$ is in $L_\alpha$, and $\Sigma$-recursion of length $\mu + 1$ gives $\Tab(\nu) = \{x \in \Lv_\nu :
\Phi(W{\restriction}(\nu+1), \nu, x, \Tab{\restriction}\nu)\}$, by $\Delta_0$-separation. (iii) Being a class
certificate is $\Delta_0$; use (i), which holds in any transitive set since $\Delta_0$ conditions are absolute, and (ii).
\end{proof}

For example, the classes $C_n$ use the clause $\Phi_n(W, \nu, x, H)$: $\nu$ is a limit above $\omega$, $p, x \in
\Lv_\nu$, and for every $c \in \Lv_\nu$ such that $\rho(c, x, p)$ is true according to $\mathrm{Tr}_\nu$, every code
$\varphi$ of a $\Pi_n$ formula in the form of Section 2.1 and every $q \in \Lv_\nu$ with $(\varphi, \langle q\rangle)
\in \mathrm{Tr}_\nu$, there is $\eta < \nu$ with $q \in \Lv_\eta$, $c \in H(\eta)$ and $(\varphi, \langle q\rangle) \in
\mathrm{Tr}_\eta$. It is $\Delta_0$. For the true certificate and $H(\eta) = \{c : \eta \in C_n(P,c)\}$ it holds
exactly when $\nu \in C_n(P,x)$, because $\nu$ and all members of the classes are limits. Lemma 4.2 uses it.

\section{Ranks}

\begin{statement}\head{Theorem 3.1} Let $\alpha > \omega$ be admissible. The ranks of the nodes in the well-founded
parts of $\Sigma_1$ presentations at $\alpha$ with parameters in $L_\alpha$ are exactly the ordinals below
$\delta_\alpha$.
\end{statement}

\begin{proof}
Let $P = (\rho, p)$ with $\rho = \exists w\, \theta$, $\theta \in \Delta_0$, and let $a$ be in the well-founded part
at $\alpha$. A \emph{history} is a finite sequence $\langle (c_1, w_1), \dots, (c_k, w_k)\rangle$ with $\theta(c_i,
c_{i-1}, p, w_i)$ for each $i$, where $c_0 = a$; its end is $c_k$, and being a history is $\Delta_0$ in $(p, a,
\omega)$. In the Kleene--Brouwer ordering of the histories, a proper extension precedes the history it extends, and
otherwise two histories are compared at the first place where they differ, by $<_L$. It is a $\Delta_1(L_\alpha)$
well-ordering, and the rank of $a$ is at most the order type of the histories before the empty one; this is the remark
after [E, Definition 3.3], and the well-foundedness and the rank bound are checked in Lean (Appendix B). Conversely, a
$\Delta_1$ well-ordering is a $\Sigma_1$ presentation in which the rank of a node is the order type of its
predecessors.
\end{proof}

Always $\delta_\alpha \le \alpha^+$: a $\Delta_1(L_\alpha)$ well-ordering is an element of $L_{\alpha+1}$, and the
admissible $L_{\alpha^+}$ computes its rank function by $\Sigma$-recursion, so its order type is below $\alpha^+$. For
countable, locally countable admissible $\alpha$, $\delta_\alpha$ is Aguilera's supremum of $\alpha$-recursive
well-orderings, since $<_L$ enumerates $L_\alpha$ in order type $\alpha$ by a $\Delta_1(L_\alpha)$ map [E, Section 8].

\begin{statement}\head{Theorem 3.2 (Gostanian, Section 2, application (viii), and Corollary 3.2)} For every countable, locally
countable admissible $\alpha < \sigma^1_1$, $\delta_\alpha = \alpha^+$. At $\alpha = \sigma^1_1$, $\delta_\alpha <
\alpha^+$.
\end{statement}

So the equality $\delta_\alpha = \alpha^+$, which makes $\delta_\alpha$ the ``degree $\alpha^+$'' of [E] and of
greatly Mahlo cardinals, holds under the hypotheses of Theorem 3.2 below $\sigma^1_1$ and fails at $\sigma^1_1$.
Theorem 3.2 says nothing about larger $\alpha$.

\section{\texorpdfstring{$\Pi_n$-reflection of degree $\delta_\alpha$}{Pi\_n-reflection of degree delta\_alpha}}

\subsection{\texorpdfstring{Upper bounds for every $n \ge 1$}{Upper bounds for every n >= 1}}

\begin{statement}\head{Lemma 4.1} Let $\mu$ be a limit ordinal above $\omega$ and $X$ a class of ordinals. Then $\mu$ is
$\Pi_1$-reflecting on $X$ if and only if $X \cap \mu$ is unbounded in $\mu$.
\end{statement}

\begin{proof}
A $\Pi_1$ sentence true in $L_\mu$ is true in every $L_\beta$, $\beta < \mu$, containing its constant, by $\Delta_0$
absoluteness, and the sentence $\forall x\, (x = x)$ with a constant $\xi < \mu$ reflects only above $\xi$.
\end{proof}

So $\mu \in C_1(P,a)$ if and only if $\mu$ is a limit above $\omega$, $p, a \in L_\mu$, and $\mu \in \Lim C_1(P,c)$ for
every predecessor $c$ of $a$ at $\mu$.

\begin{statement}\head{Lemma 4.2} Let $n \ge 1$. For each $\Sigma_1$ formula $\rho$ there is a $\Pi_{n+1}$ formula
$\Theta_n(p,a)$ such that for every admissible $\alpha > \omega$ with $p, a \in L_\alpha$, $L_\alpha \models
\Theta_n(p,a)$ if and only if $\alpha \in C_n(P,a)$; and for every $\beta$ with $p, a \in L_\beta$, $L_\beta \models
\Theta_n(p,a)$ implies $\beta \in C_n(P,a)$.
\end{statement}

\begin{proof}
\emph{Class certificates.} Apply Lemma 2.3 to the clause $\Phi_n$ given after it. So $\Mem^n_\rho(\mu, q, \varphi,
c)$, ``there is a class certificate for $\Phi_n$ through $\mu$ with $c \in \Tab(\mu)$, $q \in \Lv_\mu$ and $(\varphi,
\langle q\rangle) \in \mathrm{Tr}_\mu$'', is $\Sigma_1$, and when it holds in a transitive set, $\mu \in C_n(P,c)$, $q
\in L_\mu$ and $L_\mu \models \varphi(q)$.

\emph{The formula.}
\[ \Theta_n(p,a) = G \wedge \forall c\, \forall\varphi\, \forall q\, [\, \neg\rho(c,a,p) \vee \neg\Sat_n(\varphi,q) \vee
   \exists \mu\, \Mem^n_\rho(\mu, q, \varphi, c)\,]. \]

\emph{Its complexity.} Write $\rho = \exists i\, \rho_0$ and the last disjunct as $\exists h\, \sigma_0$, with
$\rho_0, \sigma_0$ bounded, so that $\neg\rho$ is $\forall i\, \neg\rho_0$; $\neg\Sat_n$ has $n$ blocks, the first
existential. For $n = 1$ the bracket has the prefix $\forall i\, \exists(\cdot, h)$. For $n \ge 2$, $\forall i$ joins the
second block of $\neg\Sat_n$ and $\exists h$ its first, so the bracket is $\Sigma_n$: for $n = 2$ its prefix is
$\exists (\cdot, h)\, \forall (\cdot, i)$, and for $n = 3$ it is $\exists (\cdot, h)\, \forall (\cdot, i)\, \exists$. With
$\forall c, \varphi, q$ in front and the two blocks of $G = \forall\, \exists\, \theta_G$ merged into the first two
blocks, $\Theta_n$ is $\Pi_{n+1}$.

\emph{Soundness at every $\beta$.} Let $L_\beta \models \Theta_n(p,a)$. By (C3), $\beta$ is a limit above $\omega$. Let
$c$ be a predecessor of $a$ at $\beta$ and $\varphi(q)$ a $\Pi_n$ sentence true in $L_\beta$, in the form of Section
2.1. By Section 2.1, $L_\beta \models \Sat_n(\varphi, q)$, so some $\mu < \beta$ has $\Mem^n_\rho(\mu,q,\varphi,c)$ in
$L_\beta$: $\mu \in C_n(P,c)$, $q \in L_\mu$ and $L_\mu \models \varphi(q)$. So $\beta \in C_n(P,a)$. No
admissibility of $\beta$ is used.

\emph{The converse at admissible $\alpha$.} Let $\alpha \in C_n(P,a)$. Then $L_\alpha \models G$. Let $c, \varphi, q \in
L_\alpha$ with $\rho(c,a,p)$ and $\Sat_n(\varphi,q)$ true in $L_\alpha$. Then $\varphi(q)$ is a $\Pi_n$ sentence true
in $L_\alpha$, so some $\mu < \alpha$ in $C_n(P,c)$ has $q \in L_\mu$ and $L_\mu \models \varphi(q)$, and Lemma 2.3(ii)
gives the class certificate in $L_\alpha$.
\end{proof}

\begin{statement}\head{Theorem 4.3 (upper bound)} Let $n \ge 1$. Every $\Pi_{n+1}$-reflecting ordinal $\kappa$ is
$\Pi_n$-reflecting of degree $\delta_\kappa$. In particular every admissible $\kappa > \omega$ is $\Pi_1$-reflecting of
degree $\delta_\kappa$.
\end{statement}

\begin{proof}
Let $\kappa$ be $\Pi_{n+1}$-reflecting. It is $\Pi_2$-reflecting, hence admissible and above $\omega$. Let $P$ be a
$\Sigma_1$ presentation with $p \in L_\kappa$; we show $\kappa \in C_n(P,a)$ by induction on the rank of $a$ at
$\kappa$. Let $c$ be a predecessor of $a$. By induction $\kappa \in C_n(P,c)$, so $L_\kappa \models \Theta_n(p,c)$ by
Lemma 4.2. For a $\Pi_n$ sentence $\varphi(q)$ true in $L_\kappa$, the $\Pi_{n+1}$ sentence $\varphi(q) \wedge
\Theta_n(p,c)$ holds in some $L_\beta$, $\beta < \kappa$, with $p, c, q \in L_\beta$, and $\beta \in C_n(P,c)$ by Lemma
4.2. So $\kappa$ is $\Pi_n$-reflecting on $C_n(P,c)$ for every predecessor $c$, and $\kappa \in C_n(P,a)$. The last
sentence is the case $n = 1$ (Richter--Aczel 1.8).
\end{proof}

\begin{statement}\head{Lemma 4.4} Every limit $\gamma$ of admissible ordinals is $\Pi_1$-reflecting of degree
$\delta_\gamma$.
\end{statement}

\begin{proof}
Let $\gamma \in \Lim \Ad$, $a$ in the well-founded part at $\gamma$, and $c$ a predecessor of $a$. For every admissible
$\lambda < \gamma$ with $p, c \in L_\lambda$, $c$ is in the well-founded part at $\lambda$ by (C5), so $\lambda \in
C_1(P,c)$ by Theorem 4.3. These $\lambda$ are unbounded in $\gamma$.
\end{proof}

\subsection{\texorpdfstring{Lower bounds for $n = 1$}{Lower bounds for n = 1}}

Two facts of ordinal arithmetic are used, for $\mu > 0$, both checked in Lean (Appendix B). (O1) If $\omega^{\zeta_i}$
divides $\mu$ for every $i$, so does $\omega^{\sup_i \zeta_i}$. (O2) If a set of multiples of $\omega^\zeta$ is
unbounded in $\mu$, then $\omega^{\zeta+1}$ divides $\mu$.

\begin{statement}\head{Lemma 4.5} Let $P$ be a $\Sigma_1$ presentation, $\mu > \omega$ a limit with $p, a \in L_\mu$,
$a$ in the well-founded part at $\mu$, and suppose $\omega^{1+\mathrm{rank}_\mu(a)}$ divides $\mu$. Then $\mu \in
C_1(P, a)$.
\end{statement}

\begin{proof}
Induction on $\mu$. Let $c$ be a predecessor of $a$ at $\mu$. Since $\mathrm{rank}_\mu(c) < \mathrm{rank}_\mu(a)$,
$\omega^{1+\mathrm{rank}_\mu(c)+1}$ divides $\mu$, so the positive multiples of $\omega^{1+\mathrm{rank}_\mu(c)}$ are
unbounded in $\mu$. Take such a multiple $\mu' > \omega$ above any given bound with $p, c \in L_{\mu'}$. By (C5), $c$
is in the well-founded part at $\mu'$ with $\mathrm{rank}_{\mu'}(c) \le \mathrm{rank}_\mu(c)$, so
$\omega^{1+\mathrm{rank}_{\mu'}(c)}$ divides $\mu'$ and $\mu' \in C_1(P,c)$ by the induction hypothesis. So $\mu \in
\Lim C_1(P,c)$ for every predecessor $c$.
\end{proof}

\begin{statement}\head{Corollary 4.6} Let $\lhd$ be a well-founded relation that is a set, and $a$ a node of rank
$\xi$ in it. Then $C_1(P_\lhd, a)$ is the class of limits $\mu > \omega$ with $\lhd, a \in L_\mu$ such that
$\omega^{1+\xi}$ divides $\mu$. In particular no limit $\gamma$ is in $C_1(P_\lhd,a)$ when $\xi \ge \gamma + 1$.
\end{statement}

\begin{proof}
The rank of $a$ is $\xi$ at every level containing $\lhd$, so Lemma 4.5 gives one inclusion. For the other, by
induction on $\xi$: let $\mu \in C_1(P_\lhd,a)$. If $\xi = 0$, $\omega$ divides $\mu$ because $\mu$ is a limit.
Otherwise, for each predecessor $c$, of rank $\zeta_c$, $C_1(P_\lhd,c) \cap \mu$ is unbounded in $\mu$ and its members
are divisible by $\omega^{1+\zeta_c}$; by (O2) $\omega^{1+\zeta_c+1}$ divides $\mu$, and by (O1) so does
$\omega^{1+\xi}$, since $\xi = \sup_c(\zeta_c + 1)$. For the last sentence, $\omega^{1+\xi} \ge \omega^{\gamma+1} >
\gamma$.
\end{proof}

\begin{statement}\head{Example 4.7} Lemma 4.5 is not an equivalence. Let $t = \{1\}$ and $t' = \{2\}$, and let the
predecessors of $t$ be all ordinals, of each ordinal the smaller ordinals, and of $t'$ only $t$ (a $\Delta_0$
presentation with parameter $(t, t')$). At an $\varepsilon$-number $\mu$, $\mathrm{rank}_\mu(t) = \mu$ and
$\omega^{1+\mu} = \mu$ divides $\mu$, so $\mu \in C_1(P, t)$ by Lemma 4.5. Hence every limit $\mu$ of
$\varepsilon$-numbers is in $C_1(P, t')$, although $\mathrm{rank}_\mu(t') = \mu + 1$ and $\omega^{\mu+1}$ does not
divide $\mu$.
\end{statement}

\begin{statement}\head{Theorem 4.8 (least element, $n = 1$)} No limit $\gamma$ with $\omega < \gamma <
\omega_1^{CK}$ is $\Pi_1$-reflecting of degree $\delta_\gamma$. Hence the least $\alpha$ that is $\Pi_1$-reflecting of
degree $\delta_\alpha$ is $\omega_1^{CK}$, the least $\Pi_2$-reflecting ordinal.
\end{statement}

\begin{proof}
Since $\gamma + 2 < \omega_1^{CK}$ there is a recursive well-ordering $\lhd$ of $\omega$ of order type $\gamma + 2$. It
is arithmetically definable over $L_\omega$, so $\lhd \in L_{\omega+1} \subseteq L_\gamma$. Its last element has rank
$\gamma+1$, so $\gamma$ is not in its class by Corollary 4.6. The upper bound is Theorem 4.3.
\end{proof}

\begin{statement}\head{Theorem 4.9 (characterisation, $n = 1$)} Let $\gamma > \omega$ be a limit, and let $\gamma^-$
be the supremum of $\omega$ and the admissible ordinals below $\gamma$. Suppose that $L_\gamma$ contains a
well-ordering of $\omega$ of order type at least $\gamma^-$. Then $\gamma$ is $\Pi_1$-reflecting of degree
$\delta_\gamma$ if and only if $\gamma$ is admissible or a limit of admissible ordinals.
\end{statement}

\begin{proof}
If $\gamma$ is admissible or a limit of admissibles, use Theorem 4.3 or Lemma 4.4. Suppose it is neither. Then the
admissible ordinals below $\gamma$ are bounded below $\gamma$, so $\gamma^- < \gamma$ and no admissible ordinal lies in
$(\gamma^-, \gamma]$. By Corollary 4.6 it suffices to find in $L_\gamma$ a well-founded set relation with a node of
rank $\gamma + 1$. Let $x \in L_\gamma$ be the given well-ordering, regarded as a real, and let $\omega_1^x$ be the
supremum of the order types of $x$-recursive well-orderings of $\omega$. Then $\omega_1^x > \gamma^-$, since $x$ itself
is one. By the theorem of Kripke and Platek relativised to $x$ ($\omega_1^x$ is the least $\alpha > \omega$ with
$L_\alpha[x]$ admissible; see Sacks, \emph{Higher Recursion Theory}), $L_{\omega_1^x}[x]$ is admissible, and therefore
so is $L_{\omega_1^x}$, since $\nu \mapsto L_\nu$ is $\Delta_1$ in $L_{\omega_1^x}[x]$. For $\Delta_0$-collection in
$L_{\omega_1^x}$: if every element of a set in $L_{\omega_1^x}$ has a witness in $L_{\omega_1^x}$, collection in
$L_{\omega_1^x}[x]$ gives a set of ordinal stages $\nu$ whose levels $L_\nu$ contain such witnesses, one for each
element; a single larger stage $\nu_0 < \omega_1^x$ bounds them, and $L_{\nu_0} \in L_{\omega_1^x}$ contains a witness
for every element. As $\omega_1^x$ exceeds
$\gamma^-$ and no admissible lies in $(\gamma^-, \gamma]$, $\omega_1^x > \gamma$. So there is an $x$-recursive
well-ordering $\lhd$ of $\omega$ of type $\gamma + 2$. With $x \in L_\eta$, $\eta < \gamma$, $\lhd$ is definable over
$L_\eta$, so $\lhd \in L_{\eta+1} \subseteq L_\gamma$.
\end{proof}

\begin{statement}\head{Lemma 4.10} For each $k \ge 1$, $L_{\omega_k^{CK}+1}$ contains a well-ordering of $\omega$ of
order type at least $\omega_k^{CK}$.
\end{statement}

\begin{proof}
Induction on $k$, starting from the usual ordering of $\omega$, which is in $L_{\omega+1}$ and has type
$\omega_0^{CK}$. Let $y \in L_{\omega_{k-1}^{CK}+1}$ be a well-ordering of $\omega$ of type at least
$\omega_{k-1}^{CK}$. As in the proof of Theorem 4.9, $\omega_1^y$ is admissible and above $\omega_{k-1}^{CK}$, so
$\omega_1^y \ge \omega_k^{CK}$; and every $y$-recursive well-ordering of $\omega$ lies in the admissible
$L_{\omega_k^{CK}}$ together with its rank function, so $\omega_1^y = \omega_k^{CK}$. Hence a $y$-recursive linear
ordering of $\omega$ is a well-ordering if and only if $L_{\omega_k^{CK}}$ contains an isomorphism of it onto an
ordinal, and the ordered sum of all $y$-recursive well-orderings, coded as a relation on $\omega$, is definable over
$L_{\omega_k^{CK}}$. It belongs to $L_{\omega_k^{CK}+1}$ and has type at least $\omega_k^{CK}$.
\end{proof}

\begin{statement}\head{Corollary 4.11} $\omega_\omega^{CK}$ is $\Pi_1$-reflecting of degree
$\delta_{\omega_\omega^{CK}}$ and not admissible, and every smaller $\alpha$ that is $\Pi_1$-reflecting of degree
$\delta_\alpha$ is admissible. So the least $\alpha$ that is $\Pi_1$-reflecting of degree $\delta_\alpha$ and not
$\Pi_2$-reflecting is $\omega_\omega^{CK} = \gamma_1(1)$.
\end{statement}

\begin{proof}
Lemma 4.4 gives the first statement, and $\omega_\omega^{CK}$ is not admissible because $k \mapsto \omega_k^{CK}$ is
$\Sigma_1$-definable over $L_{\omega_\omega^{CK}}$ and cofinal in it. Indeed, a certificate through a proposed value
$\nu$ verifies, by a $\Delta_0$ condition, that $\nu$ is the $k$-th admissible ordinal above $\omega$; for $j > k$ the
admissible level $L_{\omega_j^{CK}}$ contains these certificates by (C2), and these levels are cofinal in
$L_{\omega_\omega^{CK}}$. If
$L_{\omega_\omega^{CK}}$ were admissible, collection on $\omega$ would bound the values $\omega_k^{CK}$ below
$\omega_\omega^{CK}$. A non-admissible limit $\gamma <
\omega_\omega^{CK}$ has $\gamma^- = \omega$ or $\gamma^- = \omega_k^{CK}$ with $k \ge 1$, and $\gamma < (\gamma^-)^+$.
In the first case Theorem 4.8 excludes it. In the second, $L_{\gamma^- + 1} \subseteq L_\gamma$, so the hypothesis of
Theorem 4.9 holds by Lemma 4.10, and Theorem 4.9 excludes it. By Lemma 4.1, $\gamma_1(1)$, the least ordinal
$\Pi_1$-reflecting on $\Ad$, is the least limit of admissibles.
\end{proof}

\subsection{\texorpdfstring{Lower bounds for $n \ge 2$}{Lower bounds for n >= 2}}

Fix $n \ge 2$ in this subsection. We follow [E, Sections 4--7]; what changes for $n \ge 3$ is how promises are kept.

\textbf{The step beyond $[E]$, for $n = 3$.} The tree searches for a $\Pi_4$-reflecting level. Each $\Pi_4$ test $\forall u\,
\exists v\, \forall w\, \exists x\, \theta$ is either answered positively, by a smaller level where it holds, or
negatively, by a $u$; a negative answer leaves the $\Pi_3$ obligation $\forall v\, \exists w\, \forall x\, \neg\theta$.
Choosing a witness $w$ for one instance $v$ leaves the $\Pi_1$ obligation $\forall x\, \neg\theta$, whose instances are
bounded checks. Each such step removes two quantifier blocks, so for fixed $n$ only a bounded number of steps must be
unwound to find what a record has promised, and the tree stays $\Sigma_1$. In $[E]$, $n = 2$, a negative answer leaves a
$\Pi_2$ obligation, and one witness step reaches a bounded check, so promises never nest.

\textbf{What the rest of the subsection proves.} Lemma 4.14: an infinite branch yields a $\Pi_{n+1}$-reflecting level
satisfying $\varphi_0$. Lemmas 4.15 and 4.16 and Theorem 4.17: if $L_\beta$ satisfies $\varphi_0$ and the
$\Pi_2$ sentence $\chi$ (certificates, a limit above every ordinal, countable levels), a descent excludes $\beta$ from
the class of the root; if moreover no such reflecting level exists up to $\beta$, the root is in the well-founded
part, so it witnesses that $\beta$ is not $\Pi_n$-reflecting of degree $\delta_\beta$. The descent carries the sentence $G$ down, because the smaller level need
not be admissible and must still contain the certificate that makes an edge true there.

\begin{statement}\head{Lemma 4.12 (descent)} Let $P$ be a presentation and $Z$ a class of pairs $(\beta, c)$ with
$\beta$ a limit ordinal above $\omega$ and $c \in L_\beta$. Suppose that for every $(\beta, c) \in Z$ there are a
predecessor $d$ of $c$ at $\beta$ and a $\Pi_n$ sentence $\tau$, with constants in $L_\beta$ and true in $L_\beta$,
such that $(\beta', d) \in Z$ for every limit $\beta'$ with $\omega < \beta' < \beta$, $d$ and the constants of $\tau$
in $L_{\beta'}$, and $L_{\beta'} \models \tau$. Then $\beta \notin C_n(P, c)$ for every $(\beta, c) \in Z$.
\end{statement}

\begin{proof}
Otherwise let $\beta$ be least such that $(\beta, c) \in Z$ and $\beta \in C_n(P,c)$ for some $c$, and take $d$ and
$\tau$. Since $d$ is a predecessor of $c$ at $\beta$, $\beta$ is $\Pi_n$-reflecting on $C_n(P,d)$. Applying this to
$\tau$ with $d$ adjoined as a constant gives $\beta' < \beta$ in $C_n(P,d)$, hence a limit above $\omega$, with $d$ and
the constants of $\tau$ in $L_{\beta'}$ and $L_{\beta'} \models \tau$. Then $(\beta', d) \in Z$ and $\beta' \in
C_n(P,d)$, contradicting the minimality of $\beta$.
\end{proof}

This is [E, Lemma 4.1] with $\Pi_n$ in place of $\Pi_2$.

\begin{statement}\head{Definition 4.13 (the tree $T_n$)} Fix a $\Pi_n$ sentence $\varphi_0(q_0)$ in the form of
Section 2.1. A \emph{promise} is a pair $(\pi, e)$ of a code $\pi$ of a $\Pi_m$ formula in that form, $1 \le m \le n$,
and a set $e$; it asserts $\pi(e)$, and $m$ is its \emph{length}. For a $\Pi_{m+1}$ code $\psi = \forall u\, \exists
v\, \psi_1$, $\mathrm{dual}(\psi)$ is the code of the $\Pi_m$ formula obtained by putting $\forall v\,
\neg\psi_1(q,u,v)$ into the prenex form of Section 2.1 with the constant $\langle q, u\rangle$; for a $\Pi_m$
code $\pi = \forall v\, \exists w\, \pi_1$ with $m \ge 2$, $\mathrm{tail}(\pi)$ is the code of $\pi_1$ with $v, w$
moved into the constant.

\emph{Nodes.} A node of depth $k$ is a sequence of records $0, \dots, k$. Record 0 is a limit $a_0 > \omega$ with
$q_0 \in L_{a_0}$. Record $j \ge 1$ is a triple $(a_j, f_{j-1}, x_j)$ of a limit $a_j > a_{j-1}$, a surjection
$f_{j-1}$ from $\omega$ onto $L_{a_{j-1}}$ with $f_{j-1} \in L_{a_j}$, and the answer $x_j \in L_{a_j}$ to task $j-1$.
The last stage is $a_k$.

\emph{Tasks.} Fix a primitive recursive enumeration of the codes $\mathrm{Refl}(i,l)$ and $\mathrm{Keep}(j,i,l)$, $i, j,
l \in \omega$, in which each code occurs at some position $k$ with $i < k$ and $j \le k$; the other positions hold
empty tasks. Task $k$ is the \emph{next task} of a node of depth $k$ and is answered by $x_{k+1}$.

\emph{Promises.} Each record $j$ makes at most one promise $\mathrm{pr}(j)$, by recursion on $j$:
\begin{itemize}
\item[(P0)] $\mathrm{pr}(0) = (\varphi_0, q_0)$, of length $n$.
\item[(P1)] If task $j-1$ is a Refl-task for $(\psi, q)$, as defined under Answers, and $x_j = (-, u)$, then $\mathrm{pr}(j) =
  (\mathrm{dual}(\psi), \langle q, u\rangle)$, of length $n$.
\item[(P2)] If task $j-1$ is $\mathrm{Keep}(j', i, l)$, $\mathrm{pr}(j') = (\pi, e)$ has length $m \ge 3$, and $x_j =
  w$, then $\mathrm{pr}(j) = (\mathrm{tail}(\pi), \langle e, f_i(l), w\rangle)$, of length $m - 2$. Here $j' \le j-1$,
  because $\mathrm{Keep}(j',i,l)$ occurs only at positions $k \ge j'$.
\item[(P3)] Otherwise record $j$ makes no promise.
\end{itemize}

\emph{Answers.} Task $k$ is answered by $x_{k+1}$ as follows.
\begin{itemize}
\item $\mathrm{Refl}(i,l)$. If $f_i(l)$ is a pair $(\psi, q)$ with $\psi$ a $\Pi_{n+1}$ code, $\psi(q) = \forall u\,
  \exists v\, \psi_1(q,u,v)$, the answer is either $(+, b)$ with $b$ a limit, $\omega \le b < a_{k+1}$, $q \in L_b$ and
  $L_b \models \psi(q)$; or $(-, u)$ for any set $u$. The task is then a \emph{Refl-task for $(\psi,q)$}.
\item $\mathrm{Keep}(j,i,l)$. If $\mathrm{pr}(j) = (\pi, e)$ is defined, of length $m$, and $v = f_i(l)$: if $m = 1$
  and $\pi(e) = \forall v'\, \theta(e,v')$, the answer is $\emptyset$, subject to $\theta(e,v)$; if $m = 2$ and $\pi(e) =
  \forall v'\, \exists w\, \theta(e,v',w)$, the answer is some $w$ with $\theta(e,v,w)$; if $m \ge 3$, the answer is
  some set $w$, and (P2) gives the new promise.
\item Otherwise the task is empty and is answered by $\emptyset$.
\end{itemize}

\emph{Edges.} The predecessors of a node are its \emph{children}, the extensions by one record satisfying these
conditions; the root $r = \emptyset$ has the nodes of depth 0 as predecessors.
\end{statement}

\emph{The example in records ($n = 3$).} Suppose task 4 is $\mathrm{Refl}(i,l)$ with $f_i(l) = (\psi, q)$, $\psi(q) = \forall u\,
\exists v\, \forall w\, \exists x\, \theta$, and $x_5 = (-, u)$. By (P1), record 5 promises $\forall v\, \exists w\,
\forall x\, \neg\theta(q,u,v,w,x)$, of length 3. Suppose task 7 is $\mathrm{Keep}(5, i', l')$ and $v = f_{i'}(l')$. The
answer $x_8 = w$ is some set, and by (P2) record 8 promises $\forall x\, \neg\theta(q,u,v,w,x)$, with constant
$\langle\langle q,u\rangle, v, w\rangle$, of length 1. A later task $\mathrm{Keep}(8, i'', l'')$ with $x = f_{i''}(l'')$
is answered by $\emptyset$, subject to $\neg\theta(q,u,v,w,x)$. To find $\mathrm{pr}(8)$ one reads task 7, which points
to record 5, and then task 4 and $x_5$: one step through (P2).

\emph{The tree is $\Sigma_1$.} ``$s$ is a node'' and ``$d$ is a child of $s$'' assert the existence of a certificate
$W$ through the last stage according to which the conditions hold. Given $W$ they are $\Delta_0$: the levels and truth
tables are in $W$; the enumeration of tasks, recognition of codes, $\mathrm{dual}$ and $\mathrm{tail}$ are read from
$\mathrm{Tr}_\omega$ (Section 2.1); each matrix condition is read from the truth table at the last stage. For the
promises, unwinding (P0)--(P2) gives: $\mathrm{pr}(j) = (\pi, e)$ if and only if for some $t \le \lfloor (n-1)/2
\rfloor$ there are records $j = j_0 > j_1 > \cdots > j_t$ of the node and pairs $(\pi_r, e_r)$, $r \le t$, with $(\pi_0,
e_0) = (\pi, e)$, such that $(\pi_t, e_t)$ is given at $j_t$ by (P0) or (P1), and for each $r < t$ task $j_r - 1$ is
$\mathrm{Keep}(j_{r+1}, i_r, l_r)$, $\pi_{r+1}$ has length at least 3, and $(\pi_r, e_r) = (\mathrm{tail}(\pi_{r+1}),
\langle e_{r+1}, f_{i_r}(l_r), x_{j_r}\rangle)$. The bound on $t$ holds because each step through (P2) lowers the
length by 2, starting from $n$ and staying at least 1. Every $e_r$ is built by pairing from elements of $L_{a_k}$, and
$L_{a_k}$ is closed under pairing because $a_k$ is a limit, so the $e_r$ range over $\Lv_{a_k}$ and the codes over
$\Lv_\omega$. For fixed $n$ this is a finite disjunction of $\Delta_0$ formulas in the node and $W$. By (C1), $T_n$ is a
$\Sigma_1$ presentation with parameter $q_0$. At $n = 2$ the tasks $\mathrm{Keep}(0,i,l)$ play the role of the
$\mathrm{Wit}(i,l)$ tasks of [E, Section 6.1].

\begin{statement}\head{Lemma 4.14 (branches)} If the root is not in the well-founded part of $T_n$ at $\mu$, there is a
limit $\lambda \le \mu$ with $q_0 \in L_\lambda$ such that $\lambda$ is $\Pi_{n+1}$-reflecting and $L_\lambda \models
\varphi_0(q_0)$.
\end{statement}

\begin{proof}
Take a sequence of nodes of depths $0, 1, 2, \dots$, each a child of the previous one, with stages $a_0 < a_1 <
\cdots$, and let $\lambda = \sup_j a_j \le \mu$; it is a limit above $\omega$ and $L_\lambda$ is the union of the
ranges of the $f_j$.

\emph{Every promise made along the branch is true in $L_\lambda$.} By induction on the length $m$. Let $\mathrm{pr}(j)
= (\pi, e)$ and let $v \in L_\lambda$, so $v = f_i(l)$ for some $i, l$. The code $\mathrm{Keep}(j,i,l)$ occurs at a
position $k$ with $i < k$ and $j \le k$, and record $k+1$ answered it. If $m = 1$, $\theta(e,v)$ holds in $L_{a_{k+1}}$
and so in $L_\lambda$. If $m = 2$, $\theta(e,v,w)$ holds for the answer $w \in L_\lambda$. If $m \ge 3$, record $k+1$
promised $\mathrm{tail}(\pi)(\langle e,v,w\rangle)$, of length $m - 2$, which is true in $L_\lambda$ by induction; so
$\exists w\, \pi_1(e,v,w)$ holds in $L_\lambda$. As $v$ was arbitrary, $\pi(e)$ holds in $L_\lambda$. In particular
$L_\lambda \models \varphi_0(q_0)$.

\emph{Reflection.} Let $\psi(q)$ be a $\Pi_{n+1}$ sentence with $q \in L_\lambda$, in the form of Section 2.1. Then
$(\psi, q) = f_i(l)$ for some $i, l$, and $\mathrm{Refl}(i,l)$ was answered. An answer $(+, b)$ gives a limit $b <
\lambda$ with $q \in L_b$ and $L_b \models \psi(q)$. An answer $(-, u)$ made a promise $\forall v\, \neg\psi_1(q,u,v)$,
true in $L_\lambda$, so $\psi(q)$ is false there. So every $\Pi_{n+1}$ sentence true in $L_\lambda$ holds in an
earlier limit level containing its constant, and $\lambda$ is $\Pi_{n+1}$-reflecting.
\end{proof}

Let $\beta$ be a limit above $\omega$ and $s \in L_\beta$ a node. $s$ is \emph{correct at $\beta$} if every promise of
$s$, including $(\varphi_0, q_0)$, is true in $L_\beta$; $\mathrm{corr}(s)$ is the conjunction of these sentences,
which is $\Pi_n$ after padding with vacuous quantifiers and merging. A Refl-task for $(\psi,q)$ is \emph{bad at
$\beta$} if $L_\beta \models \psi(q)$ and $L_b \nvDash \psi(q)$ for every limit $b$ with $\omega \le b < \beta$ and $q
\in L_b$; other tasks are never bad. $\chi$ is the $\Pi_2$ sentence $G \wedge$ ``every ordinal is smaller than some
limit ordinal'' $\wedge$ ``every $L_\eta$ is countable'', the last conjunct saying, for every ordinal $\eta$ and through a
certificate for $L_\eta$, that there is a surjection from $\omega$ onto $L_\eta$ when $L_\eta$ is nonempty. Since $n \ge
2$, $\chi$ is $\Pi_n$. If $L_\beta \models \chi$, then $\beta$ is a limit of limit ordinals above $\omega$; every
admissible level above $\omega$ satisfies the first two conjuncts.

\begin{statement}\head{Lemma 4.15 (extension)} Suppose $L_\beta \models \chi$, the node $s \in L_\beta$ is correct at
$\beta$, and its next task is not bad at $\beta$. Then $s$ has a child $d \in L_\beta$ that is correct at $\beta$ and
is a predecessor of $s$ at $\beta$.
\end{statement}

\begin{proof}
Let $k$ be the depth of $s$ and $a_k < \beta$ its last stage. Choose the answer $x$ to task $k$ as follows. An empty
task is answered by $\emptyset$. For a Refl-task for $(\psi, q)$: if $L_\beta \nvDash \psi(q)$, some $u \in L_\beta$ has
$L_\beta \models \forall v\, \neg\psi_1(q,u,v)$, and $x = (-, u)$ makes a promise true in $L_\beta$; if $L_\beta \models
\psi(q)$, the task is not bad, so there is a limit $b$ with $\omega \le b < \beta$, $q \in L_b$ and $L_b \models
\psi(q)$, and $x = (+, b)$. For $\mathrm{Keep}(j,i,l)$ with a promise $(\pi, e)$ of length $m$, which is true in
$L_\beta$ because $s$ is correct, and $v = f_i(l) \in L_\beta$: if $m = 1$, $\theta(e,v)$ holds in $L_\beta$ and $x =
\emptyset$; if $m = 2$, some $w \in L_\beta$ has $\theta(e,v,w)$ and $x = w$; if $m \ge 3$, some $w \in L_\beta$ has
$L_\beta \models \pi_1(e,v,w)$, $x = w$, and the new promise is true in $L_\beta$.

By the last conjunct of $\chi$ there is a surjection $f \in L_\beta$ from $\omega$ onto $L_{a_k}$. The answer, $f$ and
$b$ if present are finitely many elements of $L_\beta$; take $\xi$ with $a_k < \xi < \beta$ and all of them in $L_\xi$,
and, by the second conjunct, a limit $a_{k+1}$ with $\xi < a_{k+1} < \beta$. Let $d$ extend $s$ by the record
$(a_{k+1}, f, x)$; it is in $L_\beta$. Each matrix condition holds in $L_\beta$, hence in $L_{a_{k+1}}$ by $\Delta_0$
absoluteness. By $G$ a certificate through $a_{k+1}$ is in $L_\beta$, and it is the true one by (C1), so the $\Sigma_1$
edge formula holds in $L_\beta$. $d$ is correct: its old promises are those of $s$, and its new promise was chosen
true.
\end{proof}

\begin{statement}\head{Lemma 4.16 (bad tasks)} (a) A Refl-task for $(\psi,q)$ that is bad at $\beta$ is not bad at any
limit $\beta'$ with $\omega < \beta' < \beta$ and $q \in L_{\beta'}$. (b) If a task is not bad at $\beta$, there is a
$\Pi_n$ sentence $\varsigma$ with constants in $L_\beta$, true in $L_\beta$, such that the task is not bad at any limit
$\beta' \in (\omega, \beta)$ that contains the constants of $\varsigma$ and satisfies $\varsigma$.
\end{statement}

\begin{proof}
(a) By the definition of bad at $\beta$, $\psi(q)$ is false in $L_{\beta'}$, so the task is not bad there. (b) For a
task that is not a Refl-task, take $\varsigma$ trivial. For a Refl-task for $(\psi, q)$ with $\psi(q)$ false in
$L_\beta$, take $u$ with $L_\beta \models \forall v\, \neg\psi_1(q,u,v)$ and let $\varsigma$ be this $\Pi_n$ sentence;
at $\beta'$, $\psi(q)$ is false. For a Refl-task with $\psi(q)$ true in $L_\beta$, take a limit $b < \beta$ with $q \in
L_b$ and $L_b \models \psi(q)$, and let $\varsigma$ be a trivial sentence with the constant $b$; at $\beta'$, $b <
\beta'$ shows the task is not bad.
\end{proof}

\begin{statement}\head{Theorem 4.17 (descent through the tree)} Let $q_0 \in L_\beta$, $L_\beta \models \chi$ and
$L_\beta \models \varphi_0(q_0)$. Then $\beta \notin C_n(T_n, r)$. If moreover no $\Pi_{n+1}$-reflecting $\lambda \le
\beta$ has $q_0 \in L_\lambda$ and $L_\lambda \models \varphi_0(q_0)$, then $\beta$ is not $\Pi_n$-reflecting of degree
$\delta_\beta$.
\end{statement}

\begin{proof}
Let $Z$ consist of the pairs $(\beta', r)$ with $q_0 \in L_{\beta'}$, $L_{\beta'} \models \chi \wedge
\varphi_0(q_0)$, and the pairs $(\beta', s)$ with $q_0 \in L_{\beta'}$, $L_{\beta'} \models \chi$, and $s \in
L_{\beta'}$ a node correct at $\beta'$ whose next task is not bad at $\beta'$. Every ordinal occurring in $Z$ is a
limit of limits above $\omega$. We verify the hypothesis of Lemma 4.12 for $(\beta, s) \in Z$.

\emph{The predecessor.} If $s$ is a node, let $d$ be the child given by Lemma 4.15. If $s = r$, take $\xi < \beta$ with
$q_0 \in L_\xi$ and a limit $a_0$ with $\xi < a_0 < \beta$; $d = (a_0)$ is a node of depth 0 in $L_\beta$, its only
promise $\varphi_0(q_0)$ is true in $L_\beta$, so it is correct, and its edge to $r$ holds by the certificate through
$a_0$ that $G$ supplies.

\emph{The sentence.} Let $\varsigma$ be given by Lemma 4.16(b) for the next task of $d$ if that task is not bad at
$\beta$, and trivial otherwise, and let $\tau = \chi \wedge \mathrm{corr}(d) \wedge \varsigma$, with the constant $d$
adjoined. It is true in $L_\beta$, and it is $\Pi_n$ because $\chi$ is $\Pi_2$, $n \ge 2$, and $\mathrm{corr}(d)$ and
$\varsigma$ are $\Pi_n$.

\emph{Passing down.} Let $\beta'$ be a limit with $\omega < \beta' < \beta$, $d$ and the constants of $\tau$ in
$L_{\beta'}$, and $L_{\beta'} \models \tau$. Then $L_{\beta'} \models \chi$, and $q_0 \in L_{a_0} \subseteq
L_{\beta'}$ since $d \in L_{\beta'}$. $d$ is a node at $\beta'$: its last stage is below $\beta'$, $G$ gives a
certificate through it in $L_{\beta'}$, the certificate is the true one by (C1), and the bounded conditions of the
records are those already true. The certificate comes from the conjunct $G$ of $\tau$; $\beta'$ need not be
admissible. $d$ is correct at $\beta'$ because $\mathrm{corr}(d)$ holds there. If the next task of $d$ was not bad at
$\beta$, it is not bad at $\beta'$ by Lemma 4.16(b). If it was bad at $\beta$, it is a Refl-task for a pair $(\psi, q)
= f_i(l)$ with $f_i$ in $d$; since $d \in L_{\beta'}$, $q \in L_{\beta'}$, and Lemma 4.16(a) applies. So $(\beta', d)
\in Z$.

Lemma 4.12 gives $\beta \notin C_n(T_n, s)$ for every $(\beta, s) \in Z$, in particular $\beta \notin C_n(T_n, r)$.
For the second statement: if the root were not in the well-founded part at $\beta$, Lemma 4.14 would give a $\lambda$
excluded by the hypothesis. So the root, an element of $L_\beta$, is in the well-founded part at $\beta$, the parameter
$q_0$ is in $L_\beta$, and $\beta \notin C_n(T_n, r)$; by Definition 2.2, $\beta$ is not $\Pi_n$-reflecting of degree
$\delta_\beta$.
\end{proof}

For the set relation $\{(0,1)\}$, the class $C_n$ of node 0 is the class of limits above $\omega$ containing the
parameter, so every member of $C_n(P_{\{(0,1)\}}, 1)$ is $\Pi_n$-reflecting; in particular every $\alpha$ that is
$\Pi_n$-reflecting of degree $\delta_\alpha$ is $\Pi_n$-reflecting, hence admissible.

\begin{statement}\head{Theorem 4.18 (least element, $n \ge 2$)} No $\alpha$ below $K_n$ is $\Pi_n$-reflecting of
degree $\delta_\alpha$. With Theorem 4.3: the least $\alpha$ that is $\Pi_n$-reflecting of degree $\delta_\alpha$ is
$K_n$, the least $\Pi_{n+1}$-reflecting ordinal.
\end{statement}

\begin{proof}
Let $\gamma < K_n$ be a limit above $\omega$. If $\gamma$ is not $\Pi_n$-reflecting, the set relation $\{(0,1)\}$
excludes it. Otherwise $\gamma$ is $\Pi_2$-reflecting, hence admissible. Take for $\varphi_0$ the sentence $\forall x\,
\exists y\, (y = x)$, padded to $\Pi_n$, and $q_0 = \emptyset$. $G$ and ``every ordinal is below a limit'' hold at the
admissible $\gamma$, and every $L_\eta$ with $\eta < \gamma$ is countable in $L_\gamma$ by [E, Lemma 7.2], since no
ordinal below $\gamma$ is $\Pi_{n+1}$-reflecting. So $L_\gamma \models \chi \wedge \varphi_0$, no $\lambda \le \gamma$ is
$\Pi_{n+1}$-reflecting, and Theorem 4.17 applies.
\end{proof}

\begin{statement}\head{Lemma 4.19} Every $\Pi_{n+2}$-reflecting ordinal belongs to $A_n$ and is $\Pi_n$-reflecting on
$A_n$. Such ordinals exist; in particular $\gamma_1(n)$ exists.
\end{statement}

\begin{proof}
Let $\Psi_{n+1}$ be $G \wedge \forall\psi\, \forall q\, [\neg\Sat_{n+1}(\psi,q) \vee \exists b\, \exists W\,
(\Cert(W,b) \wedge b$ is a limit above $\omega \wedge q \in \Lv_b \wedge (\psi, \langle q\rangle) \in \mathrm{Tr}_b)]$.
The last disjunct is $\Sigma_1$ and joins the first block of $\neg\Sat_{n+1}$, so $\Psi_{n+1}$ is $\Pi_{n+2}$; by (C1)
and Section 2.1, every level satisfying it is $\Pi_{n+1}$-reflecting. Let $\lambda$ be $\Pi_{n+2}$-reflecting; it is
$\Pi_{n+1}$-reflecting, so $\lambda \in A_n$, and admissible above $\omega$. For a $\Pi_{n+1}$ sentence $\psi(q)$ true
in $L_\lambda$, reflecting $\psi(q) \wedge G$ gives a target that is a limit above $\omega$, and a certificate through it
is in $L_\lambda$; so $L_\lambda \models \Psi_{n+1}$. Given a
$\Pi_n$ sentence $\varphi(q)$ true in $L_\lambda$, reflecting $\varphi(q) \wedge \Psi_{n+1}$ gives a target in $A_n$.
For existence, a countable $\alpha$ with $L_\alpha$ an elementary substructure of $L_{\omega_1}$ is $\Pi_m$-reflecting
for every $m$ [E, proof of Lemma 7.1].
\end{proof}

\begin{statement}\head{Lemma 4.20}
\begin{itemize}
\item[(a)] Let $\gamma > \omega$ be admissible, with every $L_\eta$, $\eta < \gamma$, countable in $L_\gamma$. If
  $\gamma$ is $\Pi_n$-reflecting of degree $\delta_\gamma$ and $\gamma \notin A_n$, then $\gamma$ is $\Pi_n$-reflecting
  on $A_n$.
\item[(b)] Every ordinal $\gamma'$ that is $\Pi_n$-reflecting on $A_n$ is $\Pi_n$-reflecting of degree
  $\delta_{\gamma'}$.
\item[(c)] $\gamma_1(n)$ is not $\Pi_{n+1}$-reflecting.
\end{itemize}
\end{statement}

\begin{proof}
(a) Otherwise some $\Pi_n$ sentence $\varphi_0(q_0)$ is true in $L_\gamma$ and false in $L_\lambda$ for every $\lambda
\in A_n \cap \gamma$ with $q_0 \in L_\lambda$. $L_\gamma \models \chi \wedge \varphi_0(q_0)$: $G$ and ``every ordinal is
below a limit'' hold at the admissible $\gamma$, and countability is assumed. No $\Pi_{n+1}$-reflecting $\lambda \le
\gamma$ with $q_0 \in L_\lambda$ satisfies $\varphi_0(q_0)$: $\lambda < \gamma$ is excluded by the choice of
$\varphi_0$, and $\lambda = \gamma$ because $\gamma \notin A_n$. Theorem 4.17 contradicts the hypothesis.

(b) Let $\gamma'$ be $\Pi_n$-reflecting on $A_n$; it is a limit above $\omega$. Let $P$ be a $\Sigma_1$ presentation
with $p \in L_{\gamma'}$, $a$ in the well-founded part at $\gamma'$, and $c$ a predecessor of $a$. For $\lambda \in A_n
\cap \gamma'$ with $p, c \in L_\lambda$, a descending sequence from $c$ at $\lambda$ would be one at $\gamma'$ by
(C5), so $c$ is in the well-founded part at $\lambda$, and $\lambda \in C_n(P, c)$ by Theorem 4.3. Given a $\Pi_n$
sentence true in $L_{\gamma'}$, adjoin the constants $p$ and $c$ and reflect onto $A_n$: the target lies in $C_n(P,c)$.
So $\gamma' \in C_n(P,a)$.

(c) Let $\mathrm{Ref}_{n+1}(W,\mu)$ be the bounded condition ``$\mu$ is a limit above $\omega$, and for every
$\Pi_{n+1}$ code $\psi$ and $q' \in \Lv_\mu$ with $(\psi, \langle q'\rangle) \in \mathrm{Tr}_\mu$ there is a limit $b$
with $\omega < b < \mu$, $q' \in \Lv_b$ and $(\psi, \langle q' \rangle) \in \mathrm{Tr}_b$''. For a certificate through a
limit $\mu > \omega$ it holds exactly when $\mu \in A_n$ (for the converse, reflect $\psi(q') \wedge$ ``$\omega$ exists
and every ordinal has a successor''). Let $S_n = G \wedge \forall\varphi\, \forall q\, [\neg\Sat_n(\varphi,q) \vee
\exists\mu\, \exists W\, (\Cert(W,\mu) \wedge q \in \Lv_\mu \wedge \mathrm{Ref}_{n+1}(W,\mu) \wedge (\varphi,\langle
q\rangle) \in \mathrm{Tr}_\mu)]$. As in Lemma 4.2 the bracket is $\Sigma_n$, so $S_n$ is $\Pi_{n+1}$. It holds in
$L_{\gamma_1(n)}$, since $\gamma_1(n)$ is $\Pi_n$-reflecting on $A_n$ and, being $\Pi_n$-reflecting with $n \ge 2$,
admissible, so that certificates through every smaller ordinal are present. Every level satisfying $S_n$ is
$\Pi_n$-reflecting on $A_n$, by (C1) and Section 2.1. If $\gamma_1(n)$ were $\Pi_{n+1}$-reflecting, a smaller ordinal
would satisfy $S_n$, contradicting minimality.
\end{proof}

\begin{statement}\head{Theorem 4.21 (characterisation, $n \ge 2$)} Let $\gamma > \omega$ be admissible with every
$L_\eta$, $\eta < \gamma$, countable in $L_\gamma$; this includes every admissible $\gamma$ with $\omega < \gamma \le
\gamma_1(n)$. Then $\gamma$ is $\Pi_n$-reflecting of degree $\delta_\gamma$ if and only if $\gamma \in A_n$ or
$\gamma$ is $\Pi_n$-reflecting on $A_n$. The least $\alpha$ that is $\Pi_n$-reflecting of degree $\delta_\alpha$ and
not in $A_n$ is $\gamma_1(n)$.
\end{statement}

\begin{proof}
The equivalence is Theorem 4.3 and Lemma 4.20(a), (b). For the range: if some $L_\eta$, $\eta < \gamma$, is not
countable in the admissible $L_\gamma$, then [E, Lemma 7.2] gives an ordinal below $\gamma$ that is $\Pi_m$-reflecting
for every $m$, which is $\Pi_n$-reflecting on $A_n$ by Lemma 4.19; so $\gamma_1(n) < \gamma$. By Lemma 4.20(b), (c),
$\gamma_1(n)$ is $\Pi_n$-reflecting of degree $\delta_{\gamma_1(n)}$ and not in $A_n$. If $\gamma < \gamma_1(n)$ were
such an ordinal, it would be $\Pi_n$-reflecting, hence admissible, the countability hypothesis would hold, and Lemma
4.20(a) would make it $\Pi_n$-reflecting on $A_n$, contradicting the minimality of $\gamma_1(n)$.
\end{proof}

\begin{statement}\head{Corollary 4.22} The least element of $A_n$ that is a limit of elements of $A_n$ lies strictly
between $K_n$ and $\gamma_1(n)$, and $\gamma_1(n)$ is strictly below the least element of $A_n$ that is
$\Pi_n$-reflecting on $A_n$. Every $\alpha$ below $\gamma_1(n)$ that is $\Pi_n$-reflecting of degree $\delta_\alpha$ is
$\Pi_{n+1}$-reflecting.
\end{statement}

\begin{proof}
For an ordinal $\eta < \gamma_1(n)$, reflecting a trivial sentence with the constant $\eta$ onto $A_n$ gives an element
of $A_n$ above $\eta$; so the $\Pi_2$ sentence $U = G \wedge \forall \eta\, \exists \mu\, \exists W\, [\eta$ is not an
ordinal $\vee (\Cert(W,\mu) \wedge \eta < \mu \wedge \mathrm{Ref}_{n+1}(W,\mu))]$ is true in $L_{\gamma_1(n)}$. As $n
\ge 2$ it is $\Pi_n$; reflecting it onto $A_n$ gives $\beta \in A_n \cap \gamma_1(n)$ in which, by $G$ and (C1), $A_n$
is unbounded. So the least element of $A_n$ that is a limit of $A_n$ is at most $\beta < \gamma_1(n)$, and it is above
$K_n$ since nothing in $A_n$ is below $K_n$. The least element of $A_n$ that is $\Pi_n$-reflecting on $A_n$ exists by
Lemma 4.19, is at least $\gamma_1(n)$ by minimality, and differs from it by Lemma 4.20(c). The last statement is
Theorem 4.21.
\end{proof}

\begin{statement}\head{Proposition 4.23 (linear orderings)} Let $n \ge 1$ and let $\alpha > \omega$ be admissible.
Then $\alpha$ is $\Pi_n$-reflecting of degree $\delta_\alpha$ if and only if $\alpha \in C_n(Q,x)$ for every
$\Sigma_1$ presentation $Q$ that defines a linear ordering of $L_\alpha$ at $\alpha$ and every node $x$ in the
well-founded part of that ordering.
\end{statement}

\begin{proof}
One direction is a restriction. For the other we follow [E, Proposition D], whose argument uses only that
$\Pi_n$-reflection onto a class respects inclusion of classes. Let $a$ be in the well-founded part of $P$ at $\alpha$,
and let $Q$, with parameter $(p, a, \omega)$, order the histories of the proof of Theorem 3.1 by the Kleene--Brouwer
ordering and put all other sets after them, ordered by $<_L$. $Q$ is $\Sigma_1$ through order certificates, linear at
the admissible $\alpha$, and the empty history is in its well-founded part. The point is the \emph{extension clause}:
if $t$ extends a history $s$ by one edge $(c, w)$, then ``$t$ properly extends $s$, and both are histories'' is
$\Delta_0$, so $t$ precedes $s$ at every level containing $t$, without any order certificate. Hence, by induction on
$\mu$, $\mu \in C_n(Q,s)$ implies $\mu \in C_n(P, \mathrm{end}(s))$: for a predecessor $c$ of $\mathrm{end}(s)$ at $\mu$
with witness $w \in L_\mu$, $\mu$ is $\Pi_n$-reflecting on $C_n(Q,t)$; a target $\beta$ contains $t$, so $t$ is a
history in $L_\beta$ with end $c$, and $\beta \in C_n(P,c)$ by induction. Take $s$ empty and $\mu = \alpha$.
\end{proof}

\textbf{The case $n = 1$.} Theorem 4.17 needs $n \ge 2$: the sentence $\tau$ passed down contains $\chi$, which says
that limit levels exist above every ordinal and is $\Pi_2$. A $\Pi_1$ sentence cannot say this, so the target of a
reflection could be the level $a_0 + \omega$, where the depth-0 node with stage $a_0$ has no child (a child needs a limit
stage strictly between $a_0$ and $a_0 + \omega$), and the descent stops there. Theorems 4.18, 4.21 and Lemma 4.20 also
use $n \ge 2$ to infer admissibility from $\Pi_n$-reflection, and Corollary 4.22 uses it to reflect the $\Pi_2$
sentence $U$. At $n = 1$ the order of these ordinals also changes: $\gamma_1(1) = \omega_\omega^{CK}$ lies below the
least recursively inaccessible ordinal, which is both the least element of $A_1$ that is a limit of $A_1$ and the
least element of $A_1$ that is $\Pi_1$-reflecting on $A_1$. The last assertion of Corollary 4.22 still holds at $n = 1$,
by Corollary 4.11.

\section{Recursive inaccessibility}

This section is a companion to Section 4.3, for the classes $\CA(P,a)$ of Definition 2.1, whose members are
admissible limits of the classes of the predecessors. Its distinct move is in the descent: a class only asks to be a
limit of the predecessors' classes, so no sentence is passed down. The descent carries one set $x$ to a smaller target
(Lemma 5.1), and the target is admissible, being a member of a class $\CA$, so certificates are present there without
carrying $G$. In the search tree a lower bound $\ell$ on the first stage replaces the sentence $\varphi_0$.

\subsection{Descent and membership}

\begin{statement}\head{Lemma 5.1 (descent for limit classes)} Let $P$ be a presentation and $Y$ a class of pairs
$(\beta, c)$ with $c \in L_\beta$. Suppose that for every $(\beta, c) \in Y$ there are a predecessor $d$ of $c$ at
$\beta$ and a set $x \in L_\beta$ such that $(\beta', d) \in Y$ for every $\beta' < \beta$ with $\beta' \in \CA(P,d)$
and $x \in L_{\beta'}$. Then $\beta \notin \CA(P,c)$ for every $(\beta, c) \in Y$.
\end{statement}

\begin{proof}
Otherwise let $\beta$ be least with $(\beta, c) \in Y$ and $\beta \in \CA(P,c)$, and take $d$ and $x$. Then $\beta \in
\Lim \CA(P,d)$. With $x \in L_\eta$, $\eta < \beta$, there is $\beta' \in (\eta, \beta)$ in $\CA(P,d)$; then $x \in
L_{\beta'}$, so $(\beta', d) \in Y$, contradicting the minimality of $\beta$.
\end{proof}

The set $x$ does the work of the sentence $\tau$ in Lemma 4.12; the Lean form of this lemma is in Appendix B.

\begin{statement}\head{Lemma 5.2} For each $\Sigma_1$ formula $\rho$ there is a $\Pi_2$ formula $\Theta(p,a)$ such that
for every admissible $\mu > \omega$ with $p, a \in L_\mu$, $L_\mu \models \Theta(p,a)$ if and only if $\mu \in
\CA(P,a)$.
\end{statement}

\begin{proof}
Apply Lemma 2.3 with the clause $\Phi(W, \nu, x, H)$, where $W$ is a certificate through $\nu$:
$\nu \in \Ad$ according to $W$, $p, x \in \Lv_\nu$, and for each $c \in \Lv_\nu$ with $\rho(c,x,p)$ true according to
$\mathrm{Tr}_\nu$ the set of $\eta < \nu$ with $c \in H(\eta)$ is unbounded in $\nu$. Let $\Mem(\mu, c)$ be ``there is
a class certificate through $\mu$ with $c \in \Tab(\mu)$'' ($\Sigma_1$), and let $\Theta(p,a)$ be the conjunction of
$G$ and $\forall c\, \forall y\, [\neg\rho(c,a,p) \vee \exists \mu\, (y \in L_\mu \wedge \Mem(\mu,c))]$. Since
$\neg\rho$ is $\Pi_1$ and the existential part $\Sigma_1$, $\Theta$ is $\Pi_2$ (prefix $\forall c, y, i\, \exists \mu,
h$). If $L_\beta \models \Theta$ and $\beta$ is admissible, then for every predecessor $c$ of $a$ the $\mu$ given by
class certificates are true members of $\CA(P,c)$ (Lemma 2.3(iii)), unbounded in $\beta$ (take $y$ to be each
ordinal). At admissible $\mu \in \CA(P,a)$ the class certificates through all smaller ordinals are elements of
$L_\mu$, by (C2) and Lemma 2.3(ii).
\end{proof}

\subsection{The search tree}

Fix an ordinal $\ell$, a lower bound for the first stage.

\begin{statement}\head{Definition 5.3 (the tree $S_\ell$)} $S_\ell$ has the parameter $\ell$. A node of depth $k$ is a
sequence of records $0, \dots, k$. Record 0 is an ordinal $a_0 > \max(\ell, \omega)$. Record $j \ge 1$ is $(a_j,
f_{j-1}, x_j)$: an ordinal $a_j > a_{j-1}$, a surjection $f_{j-1} \in L_{a_j}$ from $\omega$ onto $L_{a_{j-1}}$, and
the answer $x_j \in L_{a_j}$ to task $j-1$. Tasks are enumerated primitive recursively, every code $\mathrm{Refl}(i,l)$
occurring at some position $k > i$, task $k$ mentioning only $f_i$ with $i < k$; they are of one kind,
$\mathrm{Refl}(i,l)$: if $f_i(l)$ is a pair $(\psi, q)$ with $\psi = \forall u\, \exists v\, \theta$, $\theta$
bounded, the answer is $(+, b)$ with $b \in \Ad$, $b$ below the new stage, $q \in L_b$, $L_b \models \psi(q)$; or $(-,
u)$, a \emph{promise} that $\forall v\, \neg\theta(q,u,v)$. Other tasks are empty. A further node condition: every
promise $(-, u)$ of the node satisfies $\neg\theta(q,u,v)$ for all $v \in L_{a_k}$. The predecessors of a node are its
children, and the root $r = \emptyset$ has the nodes of depth 0 as predecessors.
\end{statement}

Since $x_j \in L_{a_j}$, a negative answer $(-,u)$ puts $u$ in the branch, and at every later stage the node test reads
$\neg\theta(q, u, v)$ from the truth table with all arguments present. Unlike $T_n$, this tree needs no Keep-tasks: a
promise is $\Pi_1$, and the node condition checks it on the whole of the last stage. The bound $\xi^*$ in Lemmas
5.4--5.6 and Theorem 5.7 is used in one place: in Lemma 5.6, where $L_{a_k}$ must be countable in the admissible
$L_\beta$, which [E, Lemma 7.2] gives for $\beta \le \xi^*$.

\begin{statement}\head{Lemma 5.4} $S_\ell$ is a $\Sigma_1$ presentation at every level $\mu \le \xi^*$.
\end{statement}

\begin{proof}
Given a certificate $W$ through $a_k$, all conditions are bounded, membership of $b$ in $\Ad$ being read from
$W{\restriction}(b+1)$ and the enumeration of tasks from $\mathrm{Tr}_\omega$. So ``$s$ is a node'' and ``$d$ is a
child of $s$'' are $\Sigma_1$, and by (C1) they are correct.
\end{proof}

\begin{statement}\head{Lemma 5.5 (branches)} If the root is not in the well-founded part of $S_\ell$ at $\mu \le
\xi^*$, then some recursively Mahlo ordinal lies in $(\ell, \mu]$.
\end{statement}

\begin{proof}
Take a branch with stages $a_0 < a_1 < \cdots$ and $\lambda = \sup_j a_j$. Then $\ell < \lambda \le \mu$, $\lambda$ is
a limit above $\omega$, and $L_\lambda$ is the union of the ranges of the $f_j$. For $\psi(q) = \forall u\, \exists v\,
\theta$ with $q \in L_\lambda$, the pair $(\psi, q)$ is some $f_i(l)$, and its task was answered: $(+, b)$ gives
$b \in \Ad$, $b < \lambda$, with $L_b \models \psi(q)$, and $(-, u)$ gives $\neg\theta(q,u,v)$ for every $v$ in every
$L_{a_k}$, so $L_\lambda \nvDash \psi(q)$. By block contraction $\lambda$ is $\Pi_2$-reflecting on $\Ad$, that is,
recursively Mahlo.
\end{proof}

For $\beta \in \Ad$ with $\ell < \beta \le \xi^*$ and a node $s \in L_\beta$: $s$ is \emph{correct at $\beta$} if
each of its promises holds in $L_\beta$ (these are $\Pi_1$ statements). A Refl-task for $(\psi, q)$ is \emph{bad at
$\beta$} if $L_\beta \models \psi(q)$ and $L_b \nvDash \psi(q)$ for every $b \in \Ad \cap \beta$ with $q \in L_b$.

\begin{statement}\head{Lemma 5.6 (descent through $S_\ell$)} Let $Y$ consist exactly of the pairs $(\beta, r)$ with
$\beta \in \Ad$ and $\ell < \beta \le \xi^*$, and the pairs $(\beta, s)$ at the same levels where $s \in L_\beta$ is
a node correct at $\beta$ whose next task is not bad at $\beta$. Then $\beta \notin \CA(S_\ell, r)$ for each $(\beta,
r) \in Y$.
\end{statement}

\begin{proof}
We verify the hypothesis of Lemma 5.1.

\emph{Extension.} Let $(\beta, s) \in Y$ with $s$ of depth $k$. Answer the next task: $\emptyset$ if it is empty; $(-,
u)$ with $u \in L_\beta$ and $L_\beta \models \forall v\, \neg\theta(q,u,v)$ if $\psi(q)$ is false in $L_\beta$; $(+,
b)$ with $b \in \Ad \cap \beta$, $q \in L_b$ and $L_b \models \psi(q)$ if $\psi(q)$ is true, which exists because the
task is not bad. $L_{a_k}$ is countable in $L_\beta$ ($\beta$ is admissible and $\le \xi^*$), which gives the
surjection $f_k \in L_\beta$, and there is a stage $a_{k+1} < \beta$ with the answer and $f_k$ in $L_{a_{k+1}}$. The
promises of $s$ and the new one hold in $L_\beta$, hence in $L_{a_{k+1}}$, since $\Pi_1$ sentences pass to transitive
subsets. The certificate through $a_{k+1}$ is in $L_\beta$, so the child $d$ is a predecessor of $s$ at $\beta$, and
$d$ is correct at $\beta$. For the root choose $\max(\ell, \omega) < a_0 < \beta$ and put $d = (a_0)$; by (C2) the
certificate through $a_0$ is in $L_\beta$, so $d$ is a predecessor of the root at $\beta$.

\emph{The set $x$.} Let $x$ be $u$ if the next task of $d$ is a Refl-task false at $\beta$, $b$ if it is true and
not bad, and $\emptyset$ otherwise.

\emph{Passing down.} Let $\beta' < \beta$ with $\beta' \in \CA(S_\ell, d)$ and $x \in L_{\beta'}$. Then $\beta'$ is
admissible, $\beta' > a_0 > \ell$, and $d$ is a node at $\beta'$, its certificate being in $L_{\beta'}$. It is correct
at $\beta'$, since its promises are $\Pi_1$. Its next task is not bad at $\beta'$: if it was bad at $\beta$, $\psi(q)$
is false in every $L_b$ with $b \in \Ad \cap \beta$ and $q \in L_b$, and $\beta'$ is such a $b$; if it was false at
$\beta$ with witness $u \in L_{\beta'}$, then $L_{\beta'} \models \forall v\, \neg\theta(q,u,v)$; if it was true with
$b \in L_{\beta'}$, then $b < \beta'$ shows it is not bad. So $(\beta', d) \in Y$. Lemma 5.1 gives the conclusion.
\end{proof}

\subsection{The main theorem}

\begin{statement}\head{Theorem 5.7}
\begin{itemize}
\item[(a)] Every recursively Mahlo ordinal $\mu \le \xi^*$ is recursively inaccessible of degree $\delta_\mu$.
\item[(b)] If $\gamma \le \xi^*$ is admissible and recursively inaccessible of degree $\delta_\gamma$, and $\ell <
  \gamma$, then some recursively Mahlo ordinal lies in $(\ell, \gamma]$.
\item[(c)] An admissible $\gamma \le \xi^*$ is recursively inaccessible of degree $\delta_\gamma$ if and only if
  $\gamma$ is recursively Mahlo or a limit of recursively Mahlo ordinals.
\end{itemize}
\end{statement}

\begin{proof}
(a) Let $\mu \le \xi^*$ be recursively Mahlo; it is $\Pi_2$-reflecting, hence admissible (Richter--Aczel 1.8). By
induction on the rank of $a$ at $\mu$: $\mu$ is admissible and $p, a \in L_\mu$. Let $c$ be a predecessor of $a$. By
induction $\mu \in \CA(P,c)$, so $L_\mu \models \Theta(p,c)$ by Lemma 5.2. Given $y \in L_\mu$, reflect the $\Pi_2$
sentence $\Theta(p,c) \wedge$ ``$y$ exists'' onto $\Ad$: some admissible $\beta < \mu$ with $p, c, y \in L_\beta$ has
$L_\beta \models \Theta(p,c)$, so $\beta \in \CA(P,c)$ by Lemma 5.2. Hence $\mu \in \Lim \CA(P,c)$ for every
predecessor $c$.

(b) Suppose no recursively Mahlo ordinal lies in $(\ell, \gamma]$. The pair $(\gamma, r)$ is in the class $Y$ of Lemma
5.6, so $\gamma \notin \CA(S_\ell, r)$. By Lemma 5.5 the root is in the well-founded part at $\gamma$, and by Lemma 5.4
$S_\ell$ is a $\Sigma_1$ presentation, with parameter $\ell \in L_\gamma$. So $\gamma$ is not recursively inaccessible
of degree $\delta_\gamma$.

(c) If $\gamma$ is recursively inaccessible of degree $\delta_\gamma$, then by (b) every $\ell < \gamma$ has a
recursively Mahlo ordinal in $(\ell, \gamma]$, so $\gamma$ is recursively Mahlo or a limit of recursively Mahlo
ordinals. Conversely, if $\gamma$ is recursively Mahlo use (a). If $\gamma$ is admissible and a limit of recursively
Mahlo ordinals, let $a$ be in the well-founded part at $\gamma$ and $c$ a predecessor of $a$; for recursively Mahlo
$\lambda < \gamma$ above any bound with $p, c \in L_\lambda$, $c$ is in the well-founded part at $\lambda$ by (C5), and
$\lambda \in \CA(P,c)$ by (a). So $\gamma \in \CA(P,a)$.
\end{proof}

\begin{statement}\head{Theorem 5.8} The least $\alpha$ that is recursively inaccessible of degree $\delta_\alpha$ is
the least recursively Mahlo ordinal $M$.
\end{statement}

\begin{proof}
Every member of $\CA(P,a)$ is admissible. Apply Theorem 4.3 at $K_2$ to the set relation $\in$ on the ordinal 4, in
which the node 3 has rank 3. By [E, Example 3.2(1)], $K_2 \in \Mh(3)$. Reflection of a trivial sentence onto $\Mh(2)$
gives a recursively Mahlo ordinal below $K_2$. Thus $M$ exists and $M < K_2 \le \xi^*$, so by Theorem 5.7(a) $M$ is
recursively inaccessible of degree $\delta_M$, and by Theorem 5.7(c) no admissible ordinal below $M$ is.
\end{proof}

\begin{statement}\head{Corollary 5.9} For admissible $\gamma \le \xi^*$: $\gamma$ is recursively inaccessible of degree
$\delta_\gamma$ if and only if $\gamma$ is recursively Mahlo or a limit of recursively Mahlo ordinals. The least such
ordinal that is not recursively Mahlo is the least admissible limit of recursively Mahlo ordinals.
\end{statement}

\begin{proof}
The equivalence is Theorem 5.7(c). An element that is not recursively Mahlo is admissible and a limit of $\RMc$. It
remains to see that the least admissible limit of $\RMc$ exists, is below $\xi^*$ and is not recursively Mahlo. By Theorem 4.3 and [E, Example
3.2(1)], applied to the set relation $\in$ on the ordinal 5, whose node 4 has rank 4, $K_2 \in \Mh(4)$; reflecting a
trivial sentence from $K_2$ onto $\Mh(3)$ shows that $N$ exists and $N < K_2 \le \xi^*$. At $N$, reflecting a trivial
sentence with any ordinal below $N$ as a constant onto $\RMc$ shows that $\RMc$ is cofinal in $N$; with (C2) this
makes the $\Pi_2$ sentence $G \wedge$ ``$\RMc$ is unbounded'' true in $L_N$. It reflects onto $\RMc$, giving a
recursively Mahlo limit of $\RMc$ below $N$; at that ordinal the same sentence reflects onto $\Ad$, giving an
admissible limit of $\RMc$, so the least one exists and is below $N$. If it were recursively Mahlo, the last
reflection at it would give an admissible limit of $\RMc$ below it.
\end{proof}

\begin{statement}\head{Example 5.10} A single presentation whose node has rank $\mu$ at every level $\mu$ does not
give the degree. Let $P$ be the $\Delta_0$ presentation with parameter $t = \{1\}$ in which the predecessors of $t$
are all ordinals and those of each ordinal are the smaller ordinals. At $\mu$ the node $t$ is in the well-founded part
with rank $\mu$, and $\mu \in \CA(P,t)$ if and only if $\mu$ is admissible and a limit of the class $\CA(P,\zeta)$ for
every $\zeta < \mu$: the class of $t$ is the diagonal of the classes of the ordinal nodes. Let $\mu_t$ be the least
element of $\CA(P,t)$ above $M$. Then $\mu_t$ is not recursively inaccessible of degree $\delta_{\mu_t}$. Indeed,
$\RMc$ is cofinal in $N$ (proof of Corollary 5.9), so some recursively Mahlo $\kappa$ has $M < \kappa < N$; Theorem
5.7(a) puts $\kappa$ in $\CA(P,t)$, so $\mu_t \le \kappa < \xi^*$. No recursively Mahlo ordinal lies in $(M, \mu_t)$,
since it would be in $\CA(P,t)$ by Theorem 5.7(a). Nor is $\mu_t$ recursively Mahlo: it would reflect the sentence
$\Theta(p,t)$ of Lemma 5.2, with $M$ named, onto $\Ad$, giving a member of $\CA(P,t)$ in $(M, \mu_t)$. So Theorem
5.7(b) with $\ell = M$ excludes $\mu_t$, although $\mu_t$ is in the class of $t$, whose rank at $\mu_t$ is $\mu_t$.
\end{statement}

\section{Open questions}

\begin{enumerate}
\item \textbf{$\Pi_1$, outside Theorem 4.9.} Which limits $\gamma$ that are neither admissible nor limits of
  admissibles, and for which $L_\gamma$ contains no well-ordering of $\omega$ of type at least $\gamma^-$, are
  $\Pi_1$-reflecting of degree $\delta_\gamma$?
\item \textbf{Without countability.} Theorem 4.21 assumes that every earlier level is countable in $L_\gamma$, which
  fails at large admissible $\gamma$ [E, Question 9.4]. The bound $\gamma \le \xi^*$ in Corollary 5.9 gives this
  property at each admissible stage used in the descent. What are the classes beyond these ranges? In particular, does
  the direction from left to right in Corollary 5.9 hold above $\xi^*$?
\item \textbf{Higher degrees.} Degrees such as $\delta_\alpha + 1$, $\delta_\alpha \cdot k$ and $\delta_\alpha^k$ can
  be obtained by further operations on these classes; these iterates are defined in a separate draft that is not yet
  public. What are their least elements, and does the degree
  $\delta_\alpha^{\delta_\alpha}$ for ``admissible limit of'' give the least $\Pi_3$-reflecting ordinal?
\item \textbf{Presentations over the next admissible level.} If the predecessor relation at $\mu$ may be $\Sigma_1$ over
  $L_{\mu^+}$, with $\mu$ as a parameter, what is the least recursively inaccessible ordinal in this sense, and does the
  rank of the nodes still matter?
\end{enumerate}

\section{Verification status}

\begin{center}\begin{tabular}{p{0.24\linewidth}p{0.70\linewidth}}
\toprule
Part & Status \\
\midrule
All numbered results & Written proofs, not checked by a person. The upper bounds (Theorem 4.3 and Theorem 5.7(a))
are versions of known arguments; the case $n = 2$ is $[E]$. \\
Lean & No numbered result is formalised as stated. Lean 4 with Mathlib checks abstract facts behind Lemma 5.1,
  Corollary 4.6 and Theorem 3.1 (Appendix B), with only Lean's standard axioms; Lemma 4.12 is the descent lemma of
  $[E]$, checked in Lean there. \\
References & Richter--Aczel is cited from the 1973 Oslo preprint. Theorem 3.2 is cited from Gostanian's paper, where
$\alpha$ is called good when $\delta_\alpha = \alpha^+$. The relativised Kripke--Platek theorem is cited without a checked locator. Results of $[E]$ are used as stated there. \\
\bottomrule
\end{tabular}\end{center}

\textbf{Novelty.} Not established; Section 1 compares the nearest results.

No human specialist has reviewed the proofs. The passages a specialist is most likely to question are the class
certificates of Lemma 2.3 and their use at levels that are not admissible; the nested-promise tree of Section 4.3, in
particular the bounded check of promises after Definition 4.13; and the search tree of Section 5 with Lemmas 5.2 and
5.6. Beyond correctness, the open matter is whether quantifying over every $\Sigma_1$ presentation is the notion of
degree one wants; $[E]$, Question 9.3, asks the same for its notion.

\section*{Appendix A: which result uses which}

\begin{center}\begin{tabular}{p{0.26\linewidth}p{0.68\linewidth}}
\toprule
Result & Uses \\
\midrule
Lemma 2.3 & [E, (C1), (C2)]; bounded separation; $\Sigma$-recursion \\
Theorem 3.1 & Kleene--Brouwer ordering; [E, Section 8] at the admissible level; the remark after [E, Definition 3.3] \\
Lemmas 4.1, 4.2, Theorem 4.3 & Definition 2.1, Lemma 2.3, the truth predicates of Section 2.1; [E, (C1)--(C3)];
  Richter--Aczel 1.8 \\
Lemmas 4.4, 4.5, Corollary 4.6 & Theorem 4.3, Lemma 4.1, (O1)--(O2); [E, (C5)] \\
Theorems 4.8, 4.9, Corollary 4.11 & Theorem 4.3, Lemma 4.4, Corollary 4.6, Lemma 4.10; the relativised Kripke--Platek
  theorem \\
Theorem 4.17 & Definition 4.13, Lemmas 4.12, 4.14--4.16; [E, (C1), (C3)] \\
Theorems 4.18, 4.21, Corollary 4.22 & Theorems 4.3, 4.17; Lemmas 4.19, 4.20; [E, Lemma 7.2 and proof of Lemma 7.1];
  Corollary 4.22 uses the bounded test of Lemma 4.20(c) and [E, (C2)] \\
Proposition 4.23 & Theorem 3.1, Definition 2.1; [E, Section 8] \\
Lemmas 5.1--5.6, Theorem 5.7 & Lemma 2.3; [E, (C1)--(C3), (C5)]; [E, Lemma 7.2] for the countability in Lemma 5.6 \\
Theorem 5.8, Corollary 5.9 & Theorem 5.7; Theorem 4.3 and [E, Example 3.2(1)] for the existence of $M$ and $N$; [E,
  (C2)] \\
\bottomrule
\end{tabular}\end{center}

\section*{Appendix B: the Lean statements}

The Lean file \texttt{DegreeReflection.lean} formalises three abstract statements. It covers abstract facts
only: nothing in it mentions constructible levels or definability. Ordinals are either ordinals of
Mathlib or, for the descent, a type \texttt{Pt} with a well-founded order. The file copies, unchanged, the definitions
\texttt{Refl} and \texttt{cls} of the Lean file of $[E]$ (see the appendix of $[E]$): \texttt{cls holds pred ok μ a}
says that \texttt{ok μ a} holds and that \texttt{μ} reflects on the class of every predecessor of \texttt{a} at
\texttt{μ}, where reflecting on a class means that every test true at \texttt{μ} is true at a smaller member of the
class.

\textbf{Descent.} Lemma 5.1, for arbitrary tests:

\begin{verbatim}
theorem not_cls_of_family_rel
    (hZ : ∀ δ c, Z δ c → ∃ d t, pred δ d c ∧ holds t δ ∧
      ∀ β, β < δ → cls holds pred ok β d → holds t β → Z β d)
    {δ : Pt} {c : N} (hc : Z δ c) : ¬ cls holds pred ok δ c
\end{verbatim}

To apply it, a test is a set $x$, \texttt{holds x β} means $x \in L_\beta$, and \texttt{ok μ a} means that $\mu$ is admissible
and $p, a \in L_\mu$. At a limit $\mu$, a class $Y$ has a member $\beta < \mu$ with $x \in L_\beta$ for
every $x \in L_\mu$ exactly when $Y \cap \mu$ is unbounded in $\mu$, so \texttt{cls} is the class $\CA$ of
Definition 2.1. Lemma 4.12 is [E, Lemma 4.1] with $\Pi_n$ sentences as tests, and the Lean theorem of $[E]$ places no
condition on the tests.

\textbf{Limits along a set relation.} Here \texttt{r} is a relation on a type \texttt{N} and \texttt{θ} an ordinal.
\texttt{IsSuccLimit μ} says that $\mu$ is a limit ordinal, and \texttt{IsWellFounded.rank} is Mathlib's rank in a
well-founded relation.

\begin{verbatim}
def setCls (μ : Ordinal.{u}) (a : N) : Prop :=
  cls (fun η β : Ordinal.{u} => η < β) (fun _ => r) (fun μ _ => IsSuccLimit μ ∧ θ < μ) μ a

theorem setCls_iff_rank (a : N) (μ : Ordinal.{u}) :
    setCls r θ μ a ↔ θ < μ ∧ ω ^ (1 + IsWellFounded.rank r a) ∣ μ

theorem not_setCls_of_succ_le_rank {a : N} {γ : Ordinal.{u}}
    (hγ : γ + 1 ≤ IsWellFounded.rank r a) : ¬ setCls r θ γ a
\end{verbatim}

The test \texttt{η} holds at \texttt{β} when $\eta < \beta$, so reflecting on a class says that the class is
unbounded below $\mu$. For a set relation $\lhd$ and the node $a$ under consideration: if no constructible level
contains both $\lhd$ and $a$, both classes in the equality of Corollary 4.6 are empty. Otherwise let $\theta$ be the larger of
$\omega$ and the least $\eta$ with $\lhd, a \in L_{\eta+1}$; then a limit $\mu$ is above $\omega$ with $\lhd, a \in
L_\mu$ exactly when $\theta < \mu$. The nodes below $a$ lie in the field of $\lhd$, and the bound $\theta < \mu$
imposed on them changes their classes only below $\theta$ (by induction on rank), which does not affect which $\mu > \theta$ are limits of
those classes. With Lemma 4.1, $C_1(P_\lhd, a)$ is then the class of the ordinals $\mu$ with \texttt{setCls r θ μ a},
and the two theorems are Corollary 4.6. Taking the ordinals as nodes gives the iteration of ``limit of'' over the
limits above $\theta$, with intersections at limit stages:

\begin{verbatim}
def limIter (θ ξ μ : Ordinal.{u}) : Prop :=
  setCls (fun a b : Ordinal.{u} => a < b) θ μ ξ

theorem limIter_iff_unfold (θ ξ μ : Ordinal.{u}) :
    limIter θ ξ μ ↔ IsSuccLimit μ ∧ θ < μ ∧
      ∀ ζ, ζ < ξ → ∀ η, η < μ → ∃ β, β < μ ∧ limIter θ ζ β ∧ η < β

theorem limIter_iff (θ ξ μ : Ordinal.{u}) :
    limIter θ ξ μ ↔ θ < μ ∧ ω ^ (1 + ξ) ∣ μ
\end{verbatim}

Fact (O1), in the form used, is the theorem \texttt{opow\_dvd\_of\_forall\_lt} of the file, and fact (O2) is
the theorem \texttt{opow\_succ\_dvd\_of\_unbounded}.

\textbf{Kleene--Brouwer.} Finite sequences are lists over a type \texttt{E} of entries, compared by \texttt{lt}:

\begin{verbatim}
def LeftOf (t u : List E) : Prop :=
  ∃ p x y v w, u = p ++ x :: v ∧ t = p ++ y :: w ∧ lt y x

def KB (t s : List E) : Prop :=
  (∃ w, w ≠ [] ∧ t = s ++ w) ∨ LeftOf lt t s

def KBOn (T : Set (List E)) (t s : List E) : Prop :=
  t ∈ T ∧ s ∈ T ∧ KB lt t s

theorem kbOn_wellFounded (hlt : WellFounded lt)
    (hT : ∀ s w, s ++ w ∈ T → s ∈ T)
    (hwf : WellFounded fun t s : List E => t ∈ T ∧ ∃ x, t = s ++ [x]) :
    WellFounded (KBOn lt T)

theorem rank_le_rank_kb (hlt : WellFounded lt) (hT : ∀ s w, s ++ w ∈ T → s ∈ T)
    (hstep : ∀ s x, s ++ [x] ∈ T → r (f (s ++ [x])) (f s))
    (hcover : ∀ s c, s ∈ T → r c (f s) → ∃ x, s ++ [x] ∈ T ∧ f (s ++ [x]) = c) :
    ∃ hKB : WellFounded (KBOn lt T),
      ∀ s, s ∈ T → IsWellFounded.rank r (f s) ≤ (hKB.apply s).rank
\end{verbatim}

In the proof of Theorem 3.1, \texttt{T} is the set of histories, \texttt{f s} the end of the
history \texttt{s} (so \texttt{f []} is the node $a$), \texttt{lt} the canonical well-ordering $<_L$ of the entries,
and \texttt{r} the predecessor relation restricted to the well-founded part at $\alpha$, which is well-founded.
Hypotheses \texttt{hstep} and \texttt{hcover} say that a one-step extension of a history ends at a predecessor of its
end, and that every predecessor of the end of a history is the end of such an extension; the latter holds because its
witness is in $L_\alpha$. Every history other than the empty one comes before it. What is not formal is the rest of
the argument: that the ordering is linear, so that the rank of the empty history is the order type of the histories
before it, that it is $\Delta_1$, and the bound by $\delta_\alpha$.

\section*{References}

\begin{itemize}
\item $[E]$ \href{https://waffle3z.github.io/math/greatly-mahlo-pi3.html}{The least α⁺-recursively Mahlo ordinal is the least Π₃-reflecting ordinal},
  waffle3z.github.io/math.
\item J. P. Aguilera, The order of reflection, arXiv:1906.11769 (Sections 1 and 2).
\item T. Arai, Iterating the recursively Mahlo operations, arXiv:1005.1987 (Definition 1.1, Lemma 1.2, Theorem 1.3).
\item T. Arai, Conservations of first-order reflections, \emph{Journal of Symbolic Logic} 79 (2014); arXiv:1204.0205
  (Theorem 2.2).
\item T. Arai, Mahlo classes for first-order reflections, arXiv:1910.06520 (Section 1).
\item T. Arai, Lectures on Ordinal Analysis, arXiv:2304.00246v1 (Lemma 3.33, Remark 3.34).
\item R. Gostanian, The next admissible ordinal, \emph{Annals of Mathematical Logic} 17 (1979), 171--203.
\item M. Rathjen, Proof theory of reflection, preprint, 1993 (Section 2, Theorem 4.12); published in
  \emph{Annals of Pure and Applied Logic} 68 (1994), 181--224.
\item W. Richter and P. Aczel, Inductive definitions and reflecting properties of admissible ordinals, in
  \emph{Generalized Recursion Theory}, North-Holland, 1974, 301--381 (Oslo preprint, 1973: Definition 1.7, Theorems 1.8,
  1.9, 1.10, 1.11).
\item G. E. Sacks, \emph{Higher Recursion Theory}, Springer, 1990.
\end{itemize}

\end{document}
